%%
%% The CHI Zen manual.
%%
%% Compile with LuaLaTeX (Overleaf: Menu > Settings > Compiler > LuaLaTeX,
%% TeX Live version 2026 or later). See README.md for details.
%%
%% \DocumentMetadata must come before \documentclass. It switches on a
%% tagged (accessible) PDF that aims at PDF/UA-2, plus PDF/A-4f for
%% archiving, as recommended by the acmart documentation (v2.17+).
\DocumentMetadata{
  lang          = en,
  pdfstandard   = ua-2,
  pdfstandard   = a-4f,
  tagging       = on,
  tagging-setup = {math/setup=mathml-SE},
}

\documentclass[sigconf,manuscript,screen,natbib=false]{acmart}

% Load CHIZen stuff
%
%  This file imports each of the separate .tex files in this directory -- feel free to enable/disable each separately as you'd like
%



%
%   A bunch of common includes for text formatting, features. 
%
% Accented characters: \usepackage[utf8]{inputenc} is no longer needed.
% UTF-8 has been the default input encoding of LaTeX since 2018, and
% LuaLaTeX reads UTF-8 natively.

% CJK (Chinese, Japanese, Korean) characters: \cjk{改善}
% CJKutf8 only works with pdfLaTeX, so under LuaLaTeX (the engine we use
% for tagged, accessible PDFs) CJK text is set in the Harano Aji Mincho
% font that ships with TeX Live (and therefore with Overleaf).
% Final mode (TAPS) loads neither iftex nor fontspec (acmart loads both
% under LuaLaTeX) nor CJKutf8 (not on ACM's list), so there \cjk needs
% LuaLaTeX.
\ifchizenfinal
\ifLuaTeX
  \newfontfamily\chizenCJKfont{HaranoAjiMincho-Regular.otf}
  \newcommand{\cjk}[1]{{\chizenCJKfont #1}}
\else
  \newcommand{\cjk}[1]{\PackageError{CHIZen}{\string\cjk\space needs LuaLaTeX in final mode}{%
    CJKutf8 is not on ACM's TAPS list. Compile with LuaLaTeX.}}
\fi
\else
\usepackage{iftex}
\ifLuaTeX
  \usepackage{fontspec}
  \newfontfamily\chizenCJKfont{HaranoAjiMincho-Regular.otf}
  \newcommand{\cjk}[1]{{\chizenCJKfont #1}}
\else
  \usepackage{CJKutf8}
  \newcommand{\cjk}[1]{\begin{CJK}{UTF8}{min}#1\end{CJK}}
\fi
\fi


% ACCESSIBILITY
% The old `accessibility' package is obsolete and must not be loaded.

% MARKDOWN (draft mode only: ACM's TAPS does not accept it)
%
% TODO: Explore how useful this is for general-purpose HCI papers
%       - Tables/Figures need to be a separate thing, but useful for everything else?
%
% markdown v3 loads amssymb, which clashes with the newtxmath fonts that
% acmart already loads (both define \Bbbk etc.). newtxmath provides all the
% amssymb symbols, so tell LaTeX amssymb is already there.
\ifchizenfinal\else
\makeatletter
\@ifpackageloaded{newtxmath}{\@namedef{ver@amssymb.sty}{2013/01/14 v3.01 (symbols provided by newtxmath)}}{}
\makeatother
% `notes' was called `footnotes' before markdown v3.
% LaTeX inside a markdown block: maths between dollar signs (texMathDollars),
% any other command as a raw span, `\citep{key}`{=tex} (rawAttribute), and
% % comments are dropped (texComments). These replace the `hybrid' option,
% which markdown 3.7.1 soft-deprecated. Plain text is now escaped, so write
% & as it is (not \\&); a % still starts a comment, so write \% for a percent.
\usepackage[notes,definitionLists,hashEnumerators,smartEllipses,rawAttribute,texMathDollars,texComments]{markdown}

\usepackage{paracol}
\usepackage{minted}
% fvextra (under minted) moves each code line up by 0.25pt with \vspace, so
% that background colours of adjacent lines overlap; acmart warns about any
% \vspace in the text. Listings here have no background colour, so turn it
% off. (With bgcolor=..., add bgcolorboxoverlap=0.25pt to that listing.)
\fvset{backgroundcolorboxoverlap=0pt}
% acmart's monospaced font, Inconsolata, has no italic, which minted styles
% use for comments. Say so (set them upright, as LaTeX would after a "Font
% shape undefined" warning). No spaces in the size specification: it is
% read from a macro.
\makeatletter
\begingroup\ttfamily\endgroup % declares the font family
\ifcsname\encodingdefault/\ttdefault/m/it\endcsname\else
  \edef\chizen@tmp{\noexpand\DeclareFontShape{\encodingdefault}{\ttdefault}{m}{it}{<->ssub*\ttdefault/m/n}{}}%
  \chizen@tmp
\fi
\makeatother
\fi
\ifchizenfinal
% Final mode: markdown, paracol and minted are not on ACM's TAPS list.
% Using them stops with an error; the body of the environment is skipped
% (with the comment package, which acmart loads, so verbatim text is fine).
\makeatletter
\def\chizen@textnotinfinal#1#2{%
  \PackageError{CHIZen}{#1 is not available in final mode}{%
    ACM's TAPS does not accept #1, so CHI Zen does not load it when the
    paper says \string\def\string\chizenmode{final}. #2 See the CHI Zen
    manual, Section "Camera-Ready for ACM TAPS".}}
\excludecomment{markdown}
\AddToHook{env/markdown/begin}{\chizen@textnotinfinal{markdown}{Rewrite
  the Markdown block in LaTeX.}}
\NewDocumentCommand\markdownInput{O{}m}{\chizen@textnotinfinal{markdown}{Rewrite
  #2 in LaTeX.}}
\excludecomment{minted}
\AddToHook{env/minted/begin}{\chizen@textnotinfinal{minted}{Use the
  listings package (\string\usepackage{listings} and a lstlisting
  environment) instead.}}
\NewDocumentCommand\mint{}{\chizen@textnotinfinal{minted (\string\mint)}{Use
  \string\lstinline\space from the listings package instead.}}
\NewDocumentCommand\mintinline{}{\chizen@textnotinfinal{minted (\string\mintinline)}{Use
  \string\lstinline\space from the listings package instead.}}
\NewDocumentCommand\inputminted{O{}mm}{\chizen@textnotinfinal{minted (\string\inputminted)}{Use
  \string\lstinputlisting\space from the listings package instead.}}
\NewDocumentCommand\usemintedstyle{O{}m}{}
\NewDocumentEnvironment{paracol}{m +b}{\chizen@textnotinfinal{paracol}{Use
  a table or acmart's two-column layout instead.}}{}
\makeatother
\fi
\usepackage{fontawesome5}

% This package is useful for revisions! 
% It defines \comment, which acmart's comment package defines too: keep
% comment's \comment and rename changes' one to \SCcomment.
\ifchizenfinal
% Final mode: savesym is not on ACM's TAPS list, so do the same by hand,
% and accept every change ([final]: no colours or struck-out text in the
% camera-ready version). commentmarkup=footnote keeps changes from loading
% todonotes, which loads TikZ.
\let\chizenSavedComment\comment
\let\comment\undefined
\usepackage[final,commentmarkup=footnote]{changes}
\let\SCcomment\comment
\let\comment\chizenSavedComment
\else
\usepackage{savesym}
\savesymbol{comment}
\let\comment\undefined % savesym leaves \comment as \relax, which changes v4 treats as defined
\usepackage{changes}
%\usepackage[final]{changes}
\restoresymbol{SC}{comment}
\fi


% Mark likes to add these options to his papers .... worth considering
\frenchspacing
\raggedbottom

% Lennart's acronym package
\usepackage[nolist]{acronym}

\begin{acronym}
\acro{ANOVA}{Analysis of Variance}
\acro{HCI}{Human-Computer Interaction}
\end{acronym}


% Jim's quotation stuff
%\renewenvironment{quote}{%
%  \list{}{%
%    \leftmargin0.4cm   % this is the adjusting screw
%    \rightmargin\leftmargin
%  }
%  \item\relax
%}
%{\endlist}

% Jim added this -- useful highlighting function \hl{}
%\usepackage{soul,color}



% Study Variables
% ----------------------------

% scale up the sc font slightly
\ifchizenfinal
% Final mode: letltxmacro, lipsum and relsize are not on ACM's TAPS list.
% \NewCommandCopy is built into LaTeX; \chizenrelscale does what
% \relscale does.
\makeatletter
\NewCommandCopy{\oldscshape}{\scshape}
\def\chizenrelscale#1{%
  \@tempdima=#1\dimexpr\f@size pt\relax
  \@tempdimb=#1\dimexpr\f@baselineskip\relax
  \edef\chizen@tmp{\noexpand\fontsize{\the\@tempdima}{\the\@tempdimb}}%
  \chizen@tmp\selectfont}
\renewcommand{\scshape}[1]{\oldscshape\chizenrelscale{1.2}#1}
\NewDocumentCommand\lipsum{O{}}{\PackageError{CHIZen}{\string\lipsum\space
  is not available in final mode}{Replace the placeholder text with your own.}}
\makeatother
\else
\usepackage{letltxmacro,lipsum,relsize}
\LetLtxMacro{\oldscshape}{\scshape}
\renewcommand{\scshape}[1]{\oldscshape\relscale{1.2}#1}
\fi

% an experiment factor/variable
\newcommand{\f}[1]{\mbox{{\textsc{#1}}}\xspace}
% usage: \f{block}, \f{width}, etc.
% NOTE: need \xspace package to fix spacing issues
%       \mbox prevents breaking

% a code when qualitative coding 
\newcommand{\cd}[1]{\textit{``{#1}''}\xspace}
% usage: \cd{errors}, \cd{work}


%
%   Figures -- formerly IDEA Plots -- Wraps most features from PGFPlots
%       * Handles common formatting, defines some common graph types
%       * See file for detailed documentation
%
%%%%%%%%%%%%%%%%%%%%%%%%%%%%%%%%%%%%%%%%%%%%%%%%%
%  HOW TO USE
%%%%%%%%%%%%%%%%%%%%%%%%%%%%%%%%%%%%%%%%%%%%%%%%%
%
% CHIZen/init.tex loads this file. The rules for choosing and using each
% style (and the matplotlib/ggplot equivalents) are in the CHI Zen guide:
%   plugins/chizen/skills/chizen/references/figures.md
%
% Add a graph to your paper with this template:
%
%\begin{tikzpicture}
%\begin{axis}[<STYLE>,
%	% X Axis options (xlabel, xmin, xmax, xtick={})
%	% Y Axis options (ylabel, ymin, ymax, ytick={})
%       ]
%
%    % DATA GOES IN HERE (preferably read from data/*.csv)
%
%\end{axis}
%\end{tikzpicture}
%
%
% OR, for multiple plots use this template:
%
%\begin{tikzpicture}
%\begin{groupplot}[
%   group style={
%       group size= <COL> by <ROW>,
%       x descriptions at=edge bottom,
%       y descriptions at=edge left,
%       vertical sep=0pt,
%       horizontal sep=0pt},
%  <STYLE>,
%	% X Axis options (xmin, xmax, xtick={})
%	% Y Axis options (ymin, ymax, ytick={})
%      ]
%
%    % DATA GOES IN HERE
%
%\end{groupplot}
%\end{tikzpicture}
%
%
% The following styles are defined in this library:
%
%   chizen bar          clustered bar graph (comparing conditions)
%   chizen stacked bar  stacked bar graph
%   chizen box          box plot, filled boxes with a white median (distributions)
%   chizen tufte box    Tufte's minimal box plot (a dot for the median, no box)
%   chizen tufte panel  bare-bones bar panel for Tufte-style stacked panels (groupplot)
%   chizen line         line graph (trends over time)
%   chizen Likert       diverging stacked bar for Likert responses
%   range frame         Tufte's range frame (axis lines span the data only)
% and the cycle lists (palettes) chizen lines, okabe-ito, waterloo-light,
% waterloo-dark, colorbrewer-RYB and colorbrewer-RYB-plain.
%
% NAMES. The styles used to be called `IDEA ...' (after the IDEA lab), e.g.
% `IDEA bar' and the `IDEA lines' cycle list. Those names still work as
% aliases (see the end of this file), so existing papers compile unchanged,
% but new figures should use the `chizen ...' names.
%




%%%%%%%%%%%%%%%%%%%%%%%%%%%%%%%%%%%%%%%%%%%%%%%%%
%   Common Includes
%%%%%%%%%%%%%%%%%%%%%%%%%%%%%%%%%%%%%%%%%%%%%%%%%
\usepackage{pgfplots}
\usepackage{pgfplotstable}
\usepgfplotslibrary{dateplot}
\usepgfplotslibrary{groupplots}
\usepackage{pgfgantt}

\usepackage{subcaption}

\usetikzlibrary{external}
%\tikzexternalize[prefix=tikz/] % Must be in main .tex file?

\usetikzlibrary{pgfplots.statistics}
\usetikzlibrary{matrix}

\usepackage{bbding}

\usepackage{svg}

\pgfplotsset{compat=newest}

%%%%%%%%%%%%%%%%%%%%%%%%%%%%%%%%%%%%%%%%%%%%%%%%%
%   Colour Definitions
%%%%%%%%%%%%%%%%%%%%%%%%%%%%%%%%%%%%%%%%%%%%%%%%%

\definecolor{RYB1}{RGB}{141, 211, 199}
\definecolor{RYB2}{RGB}{255, 255, 179} % NOTE: near-invisible on white (1.04:1); use colorbrewer-RYB (with outlines), not -plain
\definecolor{RYB3}{RGB}{190, 186, 218}
\definecolor{RYB4}{RGB}{251, 128, 114}
\definecolor{RYB5}{RGB}{128, 177, 211}
\definecolor{RYB6}{RGB}{253, 180, 98}
\definecolor{RYB7}{RGB}{179, 222, 105}

\pgfplotscreateplotcyclelist{colorbrewer-RYB}{
{RYB1!50!black,fill=RYB1},
{RYB2!50!black,fill=RYB2},
{RYB3!50!black,fill=RYB3},
{RYB4!50!black,fill=RYB4},
{RYB5!50!black,fill=RYB5},
{RYB6!50!black,fill=RYB6},
{RYB7!50!black,fill=RYB7},
}

\pgfplotscreateplotcyclelist{colorbrewer-RYB-plain}{
{RYB1},
{RYB2},
{RYB3},
{RYB4},
{RYB5},
{RYB6},
{RYB7},
}

%%%%%%%%%%%%%%%%%%%%
% WATERLOO PALETTES - DIGITAL
%%%%%%%%%%%%%%%%%%%%

\definecolor{AHSLight}{RGB}{0, 154, 166}
\definecolor{AHSDark}{RGB}{0, 127, 138}

\definecolor{ArtsLight}{RGB}{233, 131, 0}
\definecolor{ArtsDark}{RGB}{172, 97, 0}

\definecolor{EngineeringLight}{RGB}{204, 170, 255}
\definecolor{EngineeringDark}{RGB}{87, 6, 140}

\definecolor{EnvironmentLight}{RGB}{182, 191, 0}
\definecolor{EnvironmentDark}{RGB}{116, 120, 0}

\definecolor{MathLight}{RGB}{255, 136, 221}
\definecolor{MathDark}{RGB}{224, 36, 154}
 
\definecolor{ScienceLight}{RGB}{119, 187, 225}
\definecolor{ScienceDark}{RGB}{0, 115, 207}


\definecolor{SchoolRedLight}{RGB}{247, 119, 119}
\definecolor{SchoolRedDark}{RGB}{150, 23, 46}



\pgfplotscreateplotcyclelist{waterloo-light}{
{AHSLight!50!black,fill=AHSLight},
{ArtsLight!50!black,fill=ArtsLight},
{EngineeringLight!50!black,fill=EngineeringLight},
{EnvironmentLight!50!black,fill=EnvironmentLight},
{MathLight!50!black,fill=MathLight},
{ScienceLight!50!black,fill=ScienceLight},
{SchoolRedLight!50!black,fill=SchoolRedLight},
}

\pgfplotscreateplotcyclelist{waterloo-dark}{
{AHSDark!50!black,fill=AHSDark},
{ArtsDark!50!black,fill=ArtsDark},
{EngineeringDark!50!black,fill=EngineeringDark},
{EnvironmentDark!50!black,fill=EnvironmentDark},
{MathDark!50!black,fill=MathDark},
{ScienceDark!50!black,fill=ScienceDark},
{SchoolRedDark!50!black,fill=SchoolRedDark},
}


%%%%%%%%%%%%%%%%%%%%
% GREYSCALE / COLOUR-BLIND NOTES (relative luminance, WCAG 2 formula)
%   None of the categorical palettes above is greyscale-safe: some colours
%   have almost the same luminance and print as the same grey, e.g.
%     waterloo-light: ArtsLight .336 vs SchoolRedLight .343,
%                     MathLight .441 vs ScienceLight .449
%     waterloo-dark:  AHS, Arts, Environment, Science and Math are all .17-.19
%     colorbrewer-RYB: RYB1 .564 vs RYB6 .544; RYB2 (.960) is almost white.
%   The colours are kept as they are because published figures use them.
%   For colour-blind readers use okabe-ito below; for black-and-white print,
%   also vary marks, dash patterns or hatching, or label series directly.
%%%%%%%%%%%%%%%%%%%%

%%%%%%%%%%%%%%%%%%%%
% OKABE-ITO PALETTE - colour-blind safe (Okabe & Ito, "Color Universal Design")
%   Distinguishable with the common forms of colour blindness, but NOT in
%   greyscale (OkabeItoOrange .416 vs OkabeItoSkyBlue .405). OkabeItoYellow
%   is light on white (1.32:1), so the cycle list gives every fill a darker
%   outline, like the palettes above.
%%%%%%%%%%%%%%%%%%%%

\definecolor{OkabeItoOrange}{RGB}{230, 159, 0}
\definecolor{OkabeItoSkyBlue}{RGB}{86, 180, 233}
\definecolor{OkabeItoBluishGreen}{RGB}{0, 158, 115}
\definecolor{OkabeItoYellow}{RGB}{240, 228, 66}
\definecolor{OkabeItoBlue}{RGB}{0, 114, 178}
\definecolor{OkabeItoVermillion}{RGB}{213, 94, 0}
\definecolor{OkabeItoReddishPurple}{RGB}{204, 121, 167}
\definecolor{OkabeItoBlack}{RGB}{0, 0, 0}

\pgfplotscreateplotcyclelist{okabe-ito}{
{OkabeItoOrange!50!black,fill=OkabeItoOrange},
{OkabeItoSkyBlue!50!black,fill=OkabeItoSkyBlue},
{OkabeItoBluishGreen!50!black,fill=OkabeItoBluishGreen},
{OkabeItoYellow!50!black,fill=OkabeItoYellow},
{OkabeItoBlue!50!black,fill=OkabeItoBlue},
{OkabeItoVermillion!50!black,fill=OkabeItoVermillion},
{OkabeItoReddishPurple!50!black,fill=OkabeItoReddishPurple},
{OkabeItoBlack,fill=OkabeItoBlack},
}






%%%%%%%%%%%%%%%%%%%%
% 5-Point DIVERGENT PALETTES - From ColorBrewer2
%%%%%%%%%%%%%%%%%%%%

\definecolor{Likert5_SD}{RGB}{166,97,26}
\definecolor{Likert5_D}{RGB}{223,194,125}
\definecolor{Likert5_N}{RGB}{245,245,245}
\definecolor{Likert5_A}{RGB}{128,205,193}
\definecolor{Likert5_SA}{RGB}{1,133,113}

\definecolor{Likert5_1_SD}{RGB}{166,97,26}
\definecolor{Likert5_1_D}{RGB}{223,194,125}
\definecolor{Likert5_1_N}{RGB}{245,245,245}
\definecolor{Likert5_1_A}{RGB}{128,205,193}
\definecolor{Likert5_1_SA}{RGB}{1,133,113}

\definecolor{Likert5_2_SD}{RGB}{123,50,148}
\definecolor{Likert5_2_D}{RGB}{194,165,207}
\definecolor{Likert5_2_N}{RGB}{247,247,247}
\definecolor{Likert5_2_A}{RGB}{166,219,160}
\definecolor{Likert5_2_SA}{RGB}{0,136,55}



%%%%%%%%%%%%%%%%%%%%
% 7-Point DIVERGENT PALETTES - From ColorBrewer2
%%%%%%%%%%%%%%%%%%%%

% Purple -> Green
\definecolor{Likert7_0}{RGB}{118,42,131}
\definecolor{Likert7_1}{RGB}{175,141,195}
\definecolor{Likert7_2}{RGB}{231,212,232}
\definecolor{Likert7_3}{RGB}{229,229,229}
\definecolor{Likert7_4}{RGB}{217,240,211}
\definecolor{Likert7_5}{RGB}{127,191,123}
\definecolor{Likert7_6}{RGB}{27,120,55}

% Orange/Brown -> Blue
\definecolor{Likert7_2_0}{RGB}{140,81,10}
\definecolor{Likert7_2_1}{RGB}{216,179,101}
\definecolor{Likert7_2_2}{RGB}{246,232,195}
\definecolor{Likert7_2_3}{RGB}{229,229,229}
\definecolor{Likert7_2_4}{RGB}{199,234,229}
\definecolor{Likert7_2_5}{RGB}{90,180,172}
\definecolor{Likert7_2_6}{RGB}{1,102,94}

% Orange -> Purple
\definecolor{Likert7_3_0}{RGB}{179,88,6}
\definecolor{Likert7_3_1}{RGB}{241,163,64}
\definecolor{Likert7_3_2}{RGB}{254,224,182}
\definecolor{Likert7_3_3}{RGB}{229,229,229}
\definecolor{Likert7_3_4}{RGB}{216,218,235}
\definecolor{Likert7_3_5}{RGB}{153,142,195}
\definecolor{Likert7_3_6}{RGB}{84,39,136}

% Blue -> Red
\definecolor{Likert7_4_0}{RGB}{33,102,172}
\definecolor{Likert7_4_1}{RGB}{103,169,207}
\definecolor{Likert7_4_2}{RGB}{209,229,240}
\definecolor{Likert7_4_3}{RGB}{229,229,229}
\definecolor{Likert7_4_4}{RGB}{253,219,199}
\definecolor{Likert7_4_5}{RGB}{239,138,98}
\definecolor{Likert7_4_6}{RGB}{178,24,43}


%%%%%%%%%%%%%%%%%%%%
% 7-Point DIVERGENT PALETTES - Jim for Avatar Paper
%%%%%%%%%%%%%%%%%%%%

% Blue -> Green to match Cleric
\definecolor{Cleric_0}{RGB}{33,102,172}
\definecolor{Cleric_1}{RGB}{103,169,207}
\definecolor{Cleric_2}{RGB}{209,229,240}
\definecolor{Cleric_3}{RGB}{229,229,229}
\definecolor{Cleric_4}{RGB}{217,240,211}
\definecolor{Cleric_5}{RGB}{127,191,123}
\definecolor{Cleric_6}{RGB}{27,120,55}

% Orange -> Red
\definecolor{Monster_0}{RGB}{178,24,43}
\definecolor{Monster_1}{RGB}{239,138,98}
\definecolor{Monster_2}{RGB}{253,219,199}
\definecolor{Monster_3}{RGB}{229,229,229}
\definecolor{Monster_4}{RGB}{216,218,235}
\definecolor{Monster_5}{RGB}{153,142,195}
\definecolor{Monster_6}{RGB}{84,39,136}


\pgfplotscreateplotcyclelist{avatar}{Cleric_0, Monster_0}

%%%%%%%%%%%%%%%%%%%%
% LINE GRAPHS - default for chizen line
%   Each series differs in colour, dash pattern AND mark, so lines can be told
%   apart when printed in black and white. Colours are the Okabe-Ito palette
%   (colour-blind safe), ordered darkest first so thin lines stay visible;
%   OkabeItoYellow is left out because a thin yellow line on white is too
%   faint (1.32:1). Marks are always drawn solid. For dense data (e.g.,
%   monthly counts) add `no markers' to the axis and rely on the dash patterns.
%%%%%%%%%%%%%%%%%%%%

\pgfplotscreateplotcyclelist{chizen lines}{
{OkabeItoBlue, solid, mark=*, mark options={solid}},
{OkabeItoVermillion, dashed, mark=square*, mark options={solid}},
{OkabeItoBluishGreen, dotted, mark=triangle*, mark options={solid}},
{OkabeItoReddishPurple, dashdotted, mark=diamond*, mark options={solid}},
{OkabeItoOrange, densely dashed, mark=o, mark options={solid}},
{OkabeItoSkyBlue, densely dotted, mark=square, mark options={solid}},
{OkabeItoBlack, loosely dashed, mark=triangle, mark options={solid}},
}

%%%%%%%%%%%%%%%%%%%%%%%%%%%%%%%%%%%%%%%%%%%%%%%%%
%   Outer South Legend Hack
%%%%%%%%%%%%%%%%%%%%%%%%%%%%%%%%%%%%%%%%%%%%%%%%%
\makeatletter
\pgfplotsset{
    every axis x label/.append style={
        alias=current axis xlabel
    },
    legend pos/outer south/.style={
        /pgfplots/legend style={
            at={%
                (%
                \@ifundefined{pgf@sh@ns@current axis xlabel}%
                {xticklabel cs:0.5}%
                {current axis xlabel.south}%
                )%
            },
            anchor=north
        }
    }
}

%%%%%%%%%%%%%%%%%%%%%%%%%%%%%%%%%%%%%%%%%%%%%%%%%
%   Grid Figures Label Hack
%%%%%%%%%%%%%%%%%%%%%%%%%%%%%%%%%%%%%%%%%%%%%%%%%

\pgfplotsset{
    groupplot xlabel/.initial={},
    every groupplot x label/.style={
        at={($({\pgfplots@group@name\space c1r\pgfplots@group@rows.west}|-{\pgfplots@group@name\space c1r\pgfplots@group@rows.outer south})!0.5!({\pgfplots@group@name\space c\pgfplots@group@columns r\pgfplots@group@rows.east}|-{\pgfplots@group@name\space c\pgfplots@group@columns r\pgfplots@group@rows.outer south})$)},
        anchor=north,
    },
    groupplot ylabel/.initial={},
    every groupplot y label/.style={
            rotate=90,
        at={($({\pgfplots@group@name\space c1r1.north}-|{\pgfplots@group@name\space c1r1.outer
west})!0.5!({\pgfplots@group@name\space c1r\pgfplots@group@rows.south}-|{\pgfplots@group@name\space c1r\pgfplots@group@rows.outer west})$)},
        anchor=south
    },
    execute at end groupplot/.code={%
      \node [/pgfplots/every groupplot x label]
{\pgfkeysvalueof{/pgfplots/groupplot xlabel}};  
      \node [/pgfplots/every groupplot y label] 
{\pgfkeysvalueof{/pgfplots/groupplot ylabel}};  
    }
}

\def\endpgfplots@environment@groupplot{%
    \endpgfplots@environment@opt%
    \pgfkeys{/pgfplots/execute at end groupplot}%
    \endgroup%
}



%%%%%%%%%%%%%%%%%%%%%%%%%%%%%%%%%%%%%%%%%%%%%%%%%
%   chizen line
%%%%%%%%%%%%%%%%%%%%%%%%%%%%%%%%%%%%%%%%%%%%%%%%%
\pgfplotsset{
  chizen line/.style={
 	cycle list name = chizen lines,
        width  = \columnwidth,
        height = 6cm,
        	axis line style = very thin,
        font = \sffamily,
        % X Axis
        %major x tick style = transparent,
        xlabel style={name=xlabel},
	% Y Axis
        ymajorgrids = true,
        scaled y ticks = false,
        separate axis lines,
	axis lines* = left,
        ymajorgrids = true,
        major tick length = -5,
        %Legend
        legend style={at={(xlabel.south)},yshift=-0.5ex, xshift=-2ex, anchor=north, legend columns = 3, draw=none, column sep = 0.25cm},
        	%legend image code/.code={%
            %        \draw[#1,fill] circle (0.075cm);
            %    },
	}
}

%%%%%%%%%%%%%%%%%%%%%%%%%%%%%%%%%%%%%%%%%%%%%%%%%
%   chizen bar / A Clustered Bar Graph
%%%%%%%%%%%%%%%%%%%%%%%%%%%%%%%%%%%%%%%%%%%%%%%%%
\pgfplotsset{
  chizen bar/.style={
	ybar = 0pt,
  	cycle list name = okabe-ito,
        width  = \columnwidth,
        height = 6cm,
        	axis line style = very thin,
       	bar width=25pt,
        font = \sffamily,
        	% X Axis
        	major x tick style = transparent,
        	enlarge x limits = {abs=1.5cm},
        	% No default xlabel/ylabel: set them in each figure. The legend is
        	% anchored below the x label (node `xlabel'), so set an xlabel.
        	xlabel style={name=xlabel},
	% Y Axis
        	ymajorgrids = true,
        	scaled y ticks = false,
        	separate axis lines,
	axis lines* = left,
        	ymajorgrids = true,
        	major tick length = 0,
        %Legend
        	legend style={at={(xlabel.south)},yshift=-0.5ex, xshift=-2ex, anchor=north, legend columns = 3, draw=none, column sep = 0.25cm},
        	legend image code/.code={\draw[#1] rectangle (0.25cm,0.25cm);}
	}
}

%%%%%%%%%%%%%%%%%%%%%%%%%%%%%%%%%%%%%%%%%%%%%%%%%
%   chizen stacked bar / A Stacked Bar Graph
%%%%%%%%%%%%%%%%%%%%%%%%%%%%%%%%%%%%%%%%%%%%%%%%%
\pgfplotsset{
  chizen stacked bar/.style={
  	ybar stacked,
  	cycle list name = okabe-ito,
        width  = \columnwidth,
        height = 6cm,
        	axis line style = very thin,
       	bar width=25pt,
       	font = \sffamily,
        	% X Axis
        	major x tick style = transparent,
        	enlarge x limits = {abs=1.5cm},
        	% No default xlabel/ylabel: set them in each figure. The legend is
        	% anchored below the x label (node `xlabel'), so set an xlabel.
        	xlabel style={name=xlabel},
	% Y Axis
        	ymajorgrids = true,
        	scaled y ticks = false,
        	separate axis lines,
	axis lines* = left,
        	ymajorgrids = true,
        	major tick length = 0,
        %Legend
        	legend style={at={(xlabel.south)},yshift=-0.5ex, xshift=-2ex, anchor=north, legend columns = 3, draw=none, column sep = 0.25cm},
        	legend image code/.code={\draw[#1] rectangle (0.25cm,0.25cm);}
	}
}


%%%%%%%%%%%%%%%%%%%%%%%%%%%%%%%%%%%%%%%%%%%%%%%%%
%   chizen Likert / Divergent Bar
%%%%%%%%%%%%%%%%%%%%%%%%%%%%%%%%%%%%%%%%%%%%%%%%%
\pgfplotsset{
  chizen Likert/.style={
 	xbar stacked,
        	axis line style = very thin,
       	bar width=5pt,
        	% X Axis
        	%major x tick style = transparent,
        	% No default xlabel: set it in each figure (the legend is anchored below it).
        	xlabel style={name=xlabel},
	xtick = {},
	font = \sffamily,
	% Y Axis
	        %axis y line=none,
	        major y tick style = transparent,
        	%ymajorgrids = true,
        	scaled y ticks = false,
        	separate axis lines,
	axis lines* = left,
        	major tick length = -2,
        %Legend
        	legend style={at={(xlabel.south)},yshift=-2ex, xshift=-2ex, anchor=north, legend columns = 3, draw=none, column sep = 0.25cm},
        	legend image code/.code={\draw[#1] rectangle (0.25cm,0.25cm);}
	}
}

%%%%%%%%%%%%%%%%%%%%%%%%%%%%%%%%%%%%%%%%%%%%%%%%%
%   Tufte Stacked Panel 
%%%%%%%%%%%%%%%%%%%%%%%%%%%%%%%%%%%%%%%%%%%%%%%%%
\pgfplotsset{
  chizen tufte panel/.style={
 ybar, 
  cycle list name = okabe-ito,
	% Generic Bar Graph Options
        width = 1.1\columnwidth,
        height = .325\columnwidth,
        font = \sffamily,
	% X Axis
	enlarge x limits = {abs=.35cm},
	axis line style = very thin,
	% Y Axis
	yticklabel=\empty,
        separate axis lines,
	y axis line style= { draw opacity=0 },
	axis lines* = left,
        axis on top,
        ymajorgrids = true,
        major tick length = 0,
        grid style = {white},
	% Legend Formatting
	title style = {yshift=-.75cm,font=\sffamily\tiny}
	}
}

%%%%%%%%%%%%%%%%%%%%%%%%%%%%%%%%%%%%%%%%%%%%%%%%%
%   Tufte Range Frame
%%%%%%%%%%%%%%%%%%%%%%%%%%%%%%%%%%%%%%%%%%%%%%%%%
\pgfplotsset{
range frame/.style={
    tick align=outside,
    axis line style={opacity=0},
    after end axis/.code={
      \draw ({rel axis cs:0,0}
          -|{axis cs:\pgfplotsdataxmin,0})
        -- ({rel axis cs:0,0}
          -|{axis cs:\pgfplotsdataxmax,0});
      \draw ({rel axis cs:0,0}
          |-{axis cs:0,\pgfplotsdataymin})
        -- ({rel axis cs:0,0}
          |-{axis cs:0,\pgfplotsdataymax});
} }
}

%%%%%%%%%%%%%%%%%%%%%%%%%%%%%%%%%%%%%%%%%%%%%%%%%
%   Tufte Box Plot 
%%%%%%%%%%%%%%%%%%%%%%%%%%%%%%%%%%%%%%%%%%%%%%%%%
\pgfplotsset{
  chizen tufte box/.style={
	% Colour Options
	ybar = 0pt,
	cycle list name = okabe-ito,
        width  = \columnwidth,
        height = 6cm,
        	axis line style = very thin,
       	bar width=1pt,
       	font = \sffamily,
        	% X Axis
        	major x tick style = transparent,
        	enlarge x limits = {abs=1.5cm},
        	xlabel style={name=xlabel},
	% Y Axis
        	ymajorgrids = true,
        	scaled y ticks = false,
        	separate axis lines,
	axis lines* = left,
        	ymajorgrids = true,
        	major tick length = 0,
        %Legend
        %	legend style={at={(xlabel.south)},yshift=2ex, xshift=-2ex, anchor=north, legend columns = 3, draw=none, column sep = 0.25cm},
        %	legend image code/.code={\draw[#1] rectangle (0.25cm,0.25cm);}
 	% Boxplot specific options
	boxplot,
	boxplot/draw direction=y, 
    	%clip = false,
    	every axis plot/.style={
      	mark=o,
      	boxplot/draw direction=y,
      	boxplot/whisker extend=0,
      	boxplot/draw/box/.code={ },
	boxplot/draw/median/.code={%
	 \draw[mark size=2pt,/pgfplots/boxplot/every median/.try]
          \pgfextra
          \pgftransformshift{
            \pgfplotsboxplotpointabbox
              {\pgfplotsboxplotvalue{median}}
              {0.5}
          }
          \pgfsetfillcolor{black}
          \pgfuseplotmark{*}
          \endpgfextra
        	;
      	},
      }
  },
}


%%%%%%%%%%%%%%%%%%%%%%%%%%%%%%%%%%%%%%%%%%%%%%%%%
%   chizen box / Box Plot
%%%%%%%%%%%%%%%%%%%%%%%%%%%%%%%%%%%%%%%%%%%%%%%%%

\pgfplotsset{
    chizen box/.style={
    	% Colour Options
	ybar = 0pt,
	cycle list name = okabe-ito,
        width  = \columnwidth,
        height = 6cm,
        axis line style = very thin,
        font = \sffamily,
        	% X Axis
        	major x tick style = transparent,
        	enlarge x limits = {abs=0.75cm},
        	xlabel style={name=xlabel},
	% Y Axis
        	ymajorgrids = true,
        	scaled y ticks = false,
        	separate axis lines,
	axis lines* = left,
        	ymajorgrids = true,
        	major tick length = 0,
        %Legend
        	legend style={legend columns = 2, draw=none, column sep = 0.25cm},
        	legend image code/.code={\fill[#1] (0,0) circle (0.1cm);}
        % draw whiskers as a single line:
        boxplot/draw/whisker/.code 2 args={%
            \draw[/pgfplots/boxplot/every whisker/.try]
                (boxplot cs:##1) -- (boxplot cs:##2)
            ;
        },%
        %
        % fill the boxes:
        boxplot/every box/.style={
            fill,
        },
        % 
        % the median should be visualized as a thick white line:
        boxplot/every median/.style={
        		thick,
		fill=white, draw=white,
        },
        %boxplot/draw/median/.code={%
        %    \draw[fill=white]
        %        (boxplot cs:\pgfplotsboxplotvalue{median}) circle (3pt)
        %    ;
        %},
        % draw the average as a circle
%        boxplot/average=auto,
%        boxplot/draw/average/.code={%
%	   \draw[fill=white,draw=gray]
%            	(boxplot cs:\pgfplotsboxplotvalue{average}) circle (2pt)
%            ;
%        },
        %
        % do not clip to avoid problems with the median:
        clip=false,
        %
        boxplot/draw direction=y,
        boxplot/whisker extend=0,
        boxplot/every whisker/.style = very thick,
        %
        %
        % width of boxes:
        boxplot/box extend=0.15,
    },
    %
    %
    rshift/.style={
        xshift=\pgfkeysvalueof{/pgfplots/rshift scale},
        legend image post style={xshift=-\pgfkeysvalueof{/pgfplots/rshift scale}},
    },
    lshift/.style={
        xshift=-\pgfkeysvalueof{/pgfplots/lshift scale},
        legend image post style={xshift=\pgfkeysvalueof{/pgfplots/lshift scale}},
    },
    rshift scale/.initial=1em,
    lshift scale/.initial=1em,
}

%%%%%%%%%%%%%%%%%%%%%%%%%%%%%%%%%%%%%%%%%%%%%%%%%
%   OLD NAMES (aliases) -- keep these: existing papers and the manual use them
%%%%%%%%%%%%%%%%%%%%%%%%%%%%%%%%%%%%%%%%%%%%%%%%%
% Each `IDEA ...' style is the `chizen ...' style under its old name. Use
% the chizen names in new figures. Don't remove an alias without checking
% that nothing in the manual or the starter paper still uses it.
\pgfplotsset{
  IDEA line/.style={chizen line},
  IDEA bar/.style={chizen bar},
  IDEA stacked bar/.style={chizen stacked bar},
  IDEA Likert/.style={chizen Likert},
  IDEA tufte panel/.style={chizen tufte panel},
  IDEA tufte box/.style={chizen tufte box},
  IDEA box/.style={chizen box},
}
% A cycle list cannot be aliased by name, so `IDEA lines' repeats the
% entries of `chizen lines'. Keep the two lists identical.
\pgfplotscreateplotcyclelist{IDEA lines}{
{OkabeItoBlue, solid, mark=*, mark options={solid}},
{OkabeItoVermillion, dashed, mark=square*, mark options={solid}},
{OkabeItoBluishGreen, dotted, mark=triangle*, mark options={solid}},
{OkabeItoReddishPurple, dashdotted, mark=diamond*, mark options={solid}},
{OkabeItoOrange, densely dashed, mark=o, mark options={solid}},
{OkabeItoSkyBlue, densely dotted, mark=square, mark options={solid}},
{OkabeItoBlack, loosely dashed, mark=triangle, mark options={solid}},
}

%%%%%%%%%%%%%%%%%%%%%%%%%%%%%%%%%%%%%%%%%%%%%%%%%
% allows multiple lines for the same Y variable in stacked plots
%%%%%%%%%%%%%%%%%%%%%%%%%%%%%%%%%%%%%%%%%%%%%%%%%
\newcommand\resetstackedplots{
\makeatletter
\pgfplots@stacked@isfirstplottrue
\makeatother
}

%%%%%%%%%%%%%%%%%%%%%%%%%%%%%%%%%%%%%%%%%%%%%%%%%
% EOF
%%%%%%%%%%%%%%%%%%%%%%%%%%%%%%%%%%%%%%%%%%%%%%%%%
\makeatother

% Under LuaLaTeX, pgfplots hands some computations to Lua. That Lua backend
% currently fails on some box plots ("linear map received invalid input
% domain inf:-inf"), so use the TeX code path, as pdfLaTeX does.
\pgfplotsset{lua backend=false}

% This caches the PGF Plots stuff to speed up compile times
%\usetikzlibrary{external}
%\tikzexternalize[prefix=submission/tikz/]


%
%   Tables
%
%appendix tables
\usepackage{longtable}
\usepackage{float}
\usepackage{multicol}
\usepackage{colortbl}


% Katja's dashed lines in table stuff
% arydshln rewrites LaTeX's table code and breaks the table structure in a
% tagged PDF (rows get the wrong number of cells, which fails PDF/UA), see
% https://latex3.github.io/tagging-project/tagging-status/ . So it is only
% loaded when tagging is off; with tagging on, \cdashlinelr{2-4} draws a
% solid booktabs \cmidrule instead, and \hdashline/\cdashline are not
% available.
\makeatletter
\IfDocumentMetadataTF{%
\newcommand{\cdashlinelr}[1]{\cmidrule(lr){#1}}
}{%
\usepackage{arydshln}
\def\adl@drawiv#1#2#3{%
	\hskip.5\tabcolsep
	\xleaders#3{#2.5\@tempdimb #1{1}#2.5\@tempdimb}%
	#2\z@ plus1fil minus1fil\relax
	\hskip.5\tabcolsep}
\newcommand{\cdashlinelr}[1]{%
	\noalign{\vskip\aboverulesep
		\global\let\@dashdrawstore\adl@draw
		\global\let\adl@draw\adl@drawiv}
	\cdashline{#1}
	\noalign{\global\let\adl@draw\@dashdrawstore
		\vskip\belowrulesep}}
}
\makeatother


%
%   Stats
%
% acmart already loads etoolbox (ACM's TAPS asks papers not to load it again).
\ifchizenfinal\else
\usepackage{etoolbox}
\fi
\usepackage{siunitx}

\sisetup{
    % range-phrase = --, % make ranges use en-dash instead of "to" (e.g., ages)
    round-mode = places,
    round-precision = 1, % can change precision globally
    round-half = even
}

% Sizes
\newcommand{\by}{\,$\times$\,}
\newcommand{\size}[3][]{\SIrange[range-phrase=\by,%
                                 range-units=single]{#2}{#3}{#1}}
\newcommand{\sizen}[2][]{\SIlist[list-separator=\by,%
                                 list-final-separator=\by,%
                                 list-pair-separator=\by,%
                                %  list-units=single% % repeat (default),single,brackets
                                 ]{#2}{#1}}


% Summary Statistics
\newcommand{\sd}[2][]{$SD=\SI{#2}{#1}$}
\newcommand{\mean}[2][]{$M=\SI{#2}{#1}$}
\newcommand{\mdn}[2][]{$Mdn=\SI{#2}{#1}$}
\newcommand{\ages}[2]{Ages \numrange{#1}{#2}}
\newcommand{\gender}[2]{\num{#1} identified as #2}
\newcommand{\genderlist}[2]{\numlist{#1} identified as #2, respectively}


% Statistical Tests
\newcommand{\fprecision}{1}
\newcommand{\chisqprecision}{1}
\newcommand{\formatprob}[2]{%
    \ifdimless{#2 pt}{.001 pt}{$#1 < .001$}{%
        \ifdimless{#2 pt}{.01 pt}{$#1 < .01$}{%
            $#1 = \num[add-integer-zero=false,%
                      print-zero-integer=false,% siunitx v3: add-integer-zero alone no longer drops the 0
                      round-mode=places,%
                      round-precision=2,%
                      round-half=even]{#2}$%
        }%
    }%
}
% Effect sizes (e.g., partial eta squared) are bounded by 1: two decimals,
% no leading zero, and never "< .001" (that rule is for p-values only).
\newcommand{\formateffect}[2]{%
    $#1 = \num[add-integer-zero=false,%
               print-zero-integer=false,% siunitx v3: add-integer-zero alone no longer drops the 0
               round-mode=places,%
               round-precision=2,%
               round-half=even]{#2}$%
}
\newcommand{\p}[1]{\formatprob{p}{#1}}
\newcommand{\etasqp}[1]{\formateffect{\eta_{p}^2}{#1}}
\newcommand{\F}[3]{{$F_{#1,#2}=\num[round-mode=places,%
                                    round-precision=\fprecision]{#3}$}}
\newcommand{\owF}[2]{{$F_{#1}=\num[round-mode=places,%
                                   round-precision=\fprecision]{#2}$}}
\newcommand{\chisq}[2]{$\chi^2(#1)=\num[round-mode=places,%
                                        round-precision=\chisqprecision]{#2}$}
                                        
\newcommand{\msd}[3][]{\mean[#1]{#2}, \sd[#1]{#3}}
\newcommand{\anova}[4]{\F{#1}{#2}{#3}, \p{#4}}
\newcommand{\anovae}[5]{\F{#1}{#2}{#3}, \p{#4}, \etasqp{#5}}
\newcommand{\owanova}[4]{\owF{#1}{#2}, \p{#3}, \etasqp{#4}}
\newcommand{\chisquare}[3]{\chisq{#1}{#2}, \p{#3}}


%%
%% \BibTeX command to typeset BibTeX logo in the docs
\AtBeginDocument{%
  \providecommand\BibTeX{{%
    Bib\TeX}}}

%% Rights management information.  This information is sent to you
%% when you complete the rights form.  These commands have SAMPLE
%% values in them; it is your responsibility as an author to replace
%% the commands and values with those provided to you when you
%% complete the rights form.
% \setcopyright{acmlicensed}
% \copyrightyear{2027}
% \acmYear{2027}
% \acmDOI{XXXXXXX.XXXXXXX}


%%
%% These commands are for a JOURNAL article.
% \acmJournal{PACMHCI}
% \acmVolume{0}
% \acmNumber{0}
% \acmArticle{000}
% \acmMonth{0}

%%
%% Submission ID.
%% Use this when submitting an article to a sponsored event. You'll
%% receive a unique submission ID from the organizers
%% of the event, and this ID should be used as the parameter to this command.
%%\acmSubmissionID{123-A56-BU3}

%%
%% The majority of ACM publications use numbered citations and
%% references.  The command \citestyle{authoryear} switches to the
%% "author year" style.
%%
%% If you are preparing content for an event
%% sponsored by ACM SIGGRAPH, you must use the "author year" style of
%% citations and references.
%% Uncommenting
%% the next command will enable that style.
%%\citestyle{acmauthoryear}

% References: biblatex with ACM's styles, run with biber, as in the starter
% paper. natbib=true keeps \citet and \citep working.
\usepackage[datamodel=acmdatamodel, style=acmnumeric, backend=biber, natbib=true]{biblatex}
\addbibresource{manual/bibliography.bib}


%%
%% end of the preamble, start of the body of the document source.
\begin{document}

%%
%% The "title" command has an optional parameter,
%% allowing the author to define a "short title" to be used in page headers.
\title[Examples, Tips, and Hacks for HCI Research Papers]{CHI Zen: Examples, Tips, and Hacks for HCI Research Papers}

%%
%% The "author" command and its associated commands are used to define
%% the authors and their affiliations.
%% Of note is the shared affiliation of the first two authors, and the
%% "authornote" and "authornotemark" commands
%% used to denote shared contribution to the research.
\author{James R. Wallace}
\affiliation{%
  \institution{University of Waterloo}
  \country{Canada}
}
\email{james.wallace@uwaterloo.ca}


%\renewcommand{\shortauthors}{F. Author et al.}

%%
%% By default, the full list of authors will be used in the page
%% headers. Often, this list is too long, and will overlap
%% other information printed in the page headers. This command allows
%% the author to define a more concise list
%% of authors' names for this purpose.
%\renewcommand{\shortauthors}{Author1, Author2, & Author3}



%%
%% The abstract is a short summary of the work to be presented in the
%% article.

\begin{abstract}
\begin{figure}[H]
    \centering
    \href{https://chi2022.acm.org}{\includesvg[inkscapelatex=false, width=.45\textwidth]{figures/logo_banner_combined_theme.svg}}
    %\includegraphics[width=.32\textwidth]{figures/university-of-waterloo-vertical-logo.png} 
    \href{http://hci.uwaterloo.ca}{\includegraphics[width=.45\textwidth]{figures/uwhci-full-cs.png}}
    %\caption{Example teaser figure.}
    %\label{fig:teaser}
\end{figure}

\end{abstract}

%%
%% The code below is generated by the tool at http://dl.acm.org/ccs.cfm.
%% Please copy and paste the code instead of the example below.
%%
% \begin{CCSXML}
% <ccs2012>
% <concept>
% <concept_id>10003120.10003130.10011762</concept_id>
% <concept_desc>Human-centered computing~Empirical studies in collaborative and social computing</concept_desc>
% <concept_significance>500</concept_significance>
% </concept>
% </ccs2012>
% \end{CCSXML}

% \ccsdesc[500]{Human-centered computing~Empirical studies in collaborative and social computing}

%%
%% Keywords. The author(s) should pick words that accurately describe
%% the work being presented. Separate the keywords with commas.
%\keywords{mental health; peer-to-peer; serious game; cognitive behavioural therapy; proteus effect;}


%%
%% This command processes the author and affiliation and title
%% information and builds the first part of the formatted document.
\maketitle



%\vspace{2em}
%\begin{quote}
%    \emph{Great things are done by a series of small things brought together. --- Vincent Van Gough}
%\end{quote}

\begin{markdown}
%%%%%%%%%%%%%%%%%%%%%%%%%%%%%%%%%%%%%%%%%%%%%%%%%%%%%%%%%%%%%%%%%%%%%%%%%%%%%
%\section{Introduction}
%%%%%%%%%%%%%%%%%%%%%%%%%%%%%%%%%%%%%%%%%%%%%%%%%%%%%%%%%%%%%%%%%%%%%%%%%%%%%

# Introduction
Even the most experienced researchers are constantly looking for ways to improve their papers. And there are many different parts of a paper you can improve: from high-level concerns like experimental design, positioning the paper within the literature, and consideration of ethical issues, to sweating the details in a statistical analysis, or getting that graph *just right* to show off your earth-shattering results. But sometimes it's hard to put all the pieces together. That's where this project comes in. 
\end{markdown}

This project --- CHI Zen, a play on the words \cjk{改善} (`KaiZen') or `continuous improvement' --- describes common pitfalls and best practices based on my experience writing academic papers. That experience reflects some of my own preferences, but also a lot that I've learned from working with other HCI researchers. It also contains a lot of examples that you are free to copy, modify, or re-use as often as you'd like. The goal is to have something of an \emph{information zoo}~\citep{heer2010tour} that can make the process of creating that groundbreaking paper a little easier. 

\begin{markdown}
CHI Zen grew out of a a ‘CHI Paper Checklist’ that I developed as a resource to help both myself and students as a paper deadline approached. It encapsulates some important advice I'd collected as a PhD student and Early Career Researcher, including bigger questions like how to frame your contributions and how early you should be seeking out peer feedback, as well as fine-grained issues like common mistakes people make when citing the literature. I also included some stylistic preferences, like words and phrases to avoid in academic writing. That checklist is now provided in the Appendix.

Th bulk of document provides a sample reference for a set of templates that I developed for `PGFPlots`. Please use `PGFPlots` if you find it useful. If you do, I'd love to hear about it. I'd also love to hear from you if you have any ways to improve them. You might also prefer to use a different graphing package, and that's cool too. The examples contained in this document hopefully can provide tips or shortcuts for use in your own papers. I suspect that you can generate more or less the same graphs with your tool of choice, be it R, a Python library, or Excel. Maybe those are an even better fit for your needs; if that's the case, you should use them instead! 

Beyond that, I've slowly accumulated information about best practices, such as transparency of research and reporting of statistics, as well as ways of simplifying the (often complex) workflow of academic publishing. For instance, I've most recently added a section on change tracking, to match the recent shift in many SIG CHI conferences to the *Proceedings of the ACM* journal, which often requires authors to revise and resubmit manuscript drafts with tracked changes. I've tried to summarize what I know, point to useful packages, and provide links to additional resources where I can. 

\end{markdown}





\begin{markdown}
# How to Use CHI Zen

This package brings together examples, tips, and hacks from a variety of different sources, and regarding just about every aspect of writing a research paper. You may find it all useful, or might only want to borrow bits and pieces. It has been designed to useful from a variety of perspectives, and these are only a few of the intended use cases:  

`\subsection*{As a Template}`{=tex}

CHI Zen is written using the same `\href{https://www.acm.org/publications/authors/submissions}{ACM Template that SIG CHI Conferences use}`{=tex}. It is divided into two parts: this manual, and a new submission. The new submission is contained in the ‘submission’ folder, with placeholders for common paper sections like an Introduction, Related Work, and Discussion. These sections are pulled together in “`00-my-new-paper.tex`” in the main folder, which also pre-loads all of the `\LaTeX{}`{=tex} packages and visualization templates described here. It is intended to be a quick way to start writing your own CHI paper.

`\subsection*{As an Information Zoo}`{=tex}

Sections `\ref{sec:code}`{=tex}, `\ref{sec:tables}`{=tex}, and especially `\ref{sec:figures}`{=tex} are an *information zoo* --- they contain annotated, working examples that you can use to help you think about how to best present data in your research. The intention *is not* to define right or wrong ways to present your work, but to provide a starting point, and point to the info vis literature to explain some of the considerations you should make when implementing those graphs in your own paper.


`\subsection*{As a Quick Reference and a Checklist}`{=tex}

Sections `\ref{sec:referencing}`{=tex} and `\ref{sec:statistics}`{=tex} provide information about some common `\LaTeX{}`{=tex} questions, like referencing the literature and reporting statistics. Section `\ref{sec:transparency}`{=tex} and `\autoref{sec:checklist}`{=tex} provide some guidance on how to manage the *process* of writing a CHI paper. Section `\ref{sec:transparency}`{=tex} synthesizes some emerging best practices for transparency and replicability, and points to resources for implementing them in your own work. Section `\ref{sec:checklist}`{=tex} is an actual checklist that you can print out and use as a summative evaluation tool for your paper, preferably a week or two before the paper deadline. 


# Comments, Suggestions, Feedback

Finally, my intention in making this project available via GitHub is to allow it to grow. If you find it useful, or can suggest a way improve it, I'd love to hear about it --- drop me a line at `\href{mailto:james.wallace@uwaterloo.ca}{james.wallace@uwaterloo.ca}`{=tex}. It might also just be an important prompt to have a conversation about best practices, or an opportunity to explore different ways of thinking about or presenting research. In any case, I hope that it's an opportunity to think about how we can improve our research practices, one small step at a time. 

\end{markdown}


% \newpage 

% \begin{paracol}{2}
% \setlength{\parindent}{0em}
% \begin{CJK*}{UTF8}{gbsn}
% 其安易持， \\
% 其未兆易謀，\\
% 其脆易破，\\
% 其微易散。\\
% 為之于未有，\\
% 治之于未亂。\\
% 合抱之木，生于毫末；\\
% 九層之臺起于累土； \\
% 千里之行，始于足下。 \\
% 為者敗之，\\
% 執者失之。\\
% 聖人無為，故無敗；\\
% 無執，故無失。\\
% 民之徒事，常于幾成而敗之。\\
% 慎終如始，\\
% 則無敗事。\\
% 是以聖人欲不欲，\\
% 不貴難得之貨。\\
% 學不學，\\
% 復眾人之所過。\\
% 以輔萬物之自然，\\
% 而不敢為\\
% \flushright{-- 老子}
% \end{CJK*}

% \switchcolumn
% That settled is easily maintained, \\
% That without signs is easily conspired, \\
% That fragile is easily shattered, \\
% That insignificant is easily dispersed. \\
% Act on it before it materializes, \\
% Manage it before it becomes chaotic. \\
% A strapping tree is grown from a tiny sprout; \\
% A sky-scraping tower is built from a modest mound, \\
% A far-reaching journey begins with a small step. \\
% Those who act upon will fail, \\
% Those who hold on will lose. \\
% The master acts not, therefore fails not; \\
% Holds not on, therefore loses not. \\
% Amateurs often fail at the verge of success. \\
% Be focused in the end as in the beginning, \\
% Then there will be no failure. \\
% Hence the master desires not to be desirous, \\
% Treasures not precious possessions. \\
% Learn to be unlearned, \\
% Liberate the people of their past. \\
% Assist the myriad things in returning to their essence, \\
% And not dare act. \\
% \flushright{-- Laozi}

% \end{paracol}

% \newpage



%
% This section is written in Markdown
%

%\tableofcontents

\newpage
\begin{markdown}

# Why \LaTeX?

\LaTeX\ is far from a perfect writing tool. For one thing, it's a *typesetting environment* and not a word processor, so little things like the difference between a \` and a ' can end up being a distraction. That is, when you're writing a paper, you generally want to focus on your words, and not getting the syntax right. So why use \LaTeX?

1. Formatting is done `for free' by the SIG CHI template, and it does a great job of this for you, so you can spend your time thinking about *what* you say, not how it should look.
2. The University of Waterloo provides a Pro Overleaf account if you login with your UW userid. Overleaf is lightweight, easy-to-use, and supports collaborative writing really well. 
3. \LaTeX\ integrates well with other (free) tools like R and Python, and any referencing software you can find through the Bib\TeX package.

However, as hinted at above, \LaTeX\ also isn't perfect, and so this project embodies some examples, tips, and hacks that make our typical publication workflow a little easier. For instance, I've included some supports for all of the following tasks:

1. Writing: It supports Markdown: enabling you to focus on *content*, and the ACM templates to handle all formatting

2. Collaborative Editing: Sections are separated out into files, so that authors can edit each one without interrupting one another on Overleaf, without the need to email drafts back and forth between authors, and with features like version control, and automatic backups. 

3. Illustrating: Several commonly-used \LaTeX\ packages are automatically included, for tasks like creating figures and tables, documenting code, and referencing

4. Revising: The \texttt{changes} package helps authors to show revisions to their work when working within the ACM's Revise \\& Resubmit model

5. Many (optional) in-text shortcuts are also included, for tasks like reporting statistics, defining colour palettes, and including alternative languages. They are there if you need them, and easy to ignore if you don't. 



## Markdown
A common complaint about \LaTeX\ is that it's hard to read and write, and that the learning curve is *really* steep. For this reason, we're following the general philosophy of letting \LaTeX\ handle the stuff it's good at --- formatting, math, figures, references --- and writing the rest in Markdown. Markdown is extremely easy to read and write. In fact, you're probably familiar with much of its syntax already: 

>    "The overriding design goal for Markdown’s formatting syntax is to make it as readable as possible. The idea is that a Markdown-formatted document should be publishable as-is, as plain text, without looking like it’s been marked up with tags or formatting instructions. While Markdown’s syntax has been influenced by several existing text-to-HTML filters, the single biggest source of inspiration for Markdown’s syntax is the format of plain text email"
> --- John Gruber \citep{gruber2006}.

However, you should be aware of some its limitations --- especially for some of the more technical aspects of a scientific paper. It doesn't always play well with tables, and while you can include images, Markdown doesn't necessarily offer some of the nice features that \LaTeX\ does. Keep reading for more information on that. 

**This template has been configured to allow *both* Markdown and \LaTeX\ syntax in the same file**, through the \texttt{markdown} package's \texttt{hybrid} option. You can even combine them in the same *sentence*, and for example write paragraphs in Markdown that include BibTeX citations.



## PGFPlots and My Templates 
As as I was completing my PhD, I decided to spend some time trying to learn to make better figures since this was something I expected to spend a lot of time doing over the rest of my career. Some of the things I'd hoped to improve on were using scalable graphics, having better control over the width/height of figures, and being able to use the data itself to create the images (better fit with a typical academic workflow). I also prefer a minimalist feel, and wanted my papers to reflect that (And even if you prefer something different, this is a good starting point, you can always make things more complex where needed). 

So after doing some research, I found that \texttt{PGFPlots} was a good fit for my needs. But it wasn't perfect, and in particular can be complicated to get right. So in the interest of saving the time of re-inventing the wheel, and (hopefully) enabling others to use my work for their own papers, I decided to create a set of templates that would encapsulate the common types of graphs I use in HCI research. I called these templates \texttt{IDEA Plots} after my lab at the time. As I developed the templates, they grew to include style choices, colour palettes, and some common shortcuts for layout. I expect they'll continue to evolve for a while to come.  

All of the examples in this document use those templates, which are now wrapped into CHI Zen, with some in-place customization based on the needs of the graph. 



## Tracking Changes

Since SIG CHI has begun transitioning its conferences to the *Proceedings of the ACM* model, most involve some form of Revise \\& Resubmit. That is, you should expect to take reviewer feedback into consideration and submit an updated version of your paper for a second round of review. When doing so, conferences typically ask for a marked-up version of the paper that shows where changes have been made.  To track changes in \LaTeX, we use the \texttt{changes} package. It provides a few simple functions to track anything that you've \texttt{added()}, \texttt{deleted()}, or \texttt{replaced} in the paper. 



\end{markdown}

% Written in LaTeX rather than Markdown because it shows LaTeX commands.
\subsection{Accessibility}
\label{sec:accessibility}

Accessibility used to be a perennial to-do item for \LaTeX\ papers: the \texttt{accessibility} package never worked with this template, and the SIGCHI guide had no \LaTeX\ advice beyond figure descriptions. That has changed. The \LaTeX\ Project's tagging project now produces tagged PDF, and \texttt{acmart} supports it (since v2.17). CHI Zen turns it on for you.

\paragraph{What CHI asks for.} The CHI 2027 publication formats page states that ``The CHI conference committee requests that authors ensure basic accessibility of their submissions'', and points to the SIGCHI \href{https://sigchi.org/resources/guides-for-authors/accessibility/}{Accessibility Guide for Authors} (last updated 2024-08-15). For the camera-ready version, CHI 2025 told authors that ``An external provider will tag TAPS-produced PDFs for accessibility; however, authors must provide meaningful alt text for figures in their source files and on PCS.'' The SIGCHI guide's rules are:
\begin{itemize}
    \item ``Mark up content such as headings and lists'': use \texttt{\textbackslash section}, \texttt{itemize}, \texttt{enumerate}, \texttt{figure} and \texttt{table}, not bold text or manual numbering.
    \item ``Provide a text description for all figures.'' In \LaTeX, ``you can use the \texttt{\textbackslash Description\{...\}} tag to add an image description.'' ``Figure descriptions are different to figure captions'': they ``are an alternative to seeing the figure'' and should not ``simply repeat the caption.'' The SIGACCESS \href{https://www.sigaccess.org/welcome-to-sigaccess/resources/describing-figures/}{Guide to Describing Figures} explains how to write them.
    \item ``Don't rely only on colour'': in figures and legends, ``use different shapes and patterns to provide another way to visually distinguish elements.''
    \item ``Create every table as a real table, not an image, and indicate which cells are headers.''
    \item ``Create every equation as a marked-up equation, not an image of an equation.''
    \item ``Set the metadata of your document'': its title and language.
\end{itemize}

\paragraph{What CHI Zen does.} The first line of \texttt{00-my-new-paper.tex} is a \texttt{\textbackslash DocumentMetadata} block that asks for a tagged PDF meeting PDF/UA-2 (the accessibility standard) and PDF/A-4f (the archival standard), in English. Compile with LuaLaTeX on TeX Live 2026 or later (on Overleaf: Menu, Settings, Compiler). The PDF then gets its headings, lists, tables, links, title and language tagged automatically, and the title you give \texttt{\textbackslash title} becomes the document title that screen readers announce. Note the difference: PDF/UA is about \emph{access}, PDF/A about \emph{preservation}; a PDF/A file is not accessible unless it is also tagged.

\paragraph{What you still have to do.}
\begin{enumerate}
    \item \textbf{Describe every figure.} Put \texttt{\textbackslash Description\{...\}} as the first line of each \texttt{figure}, \texttt{figure*} and \texttt{subfigure}. \texttt{acmart} sends it to TAPS for the HTML version; CHI Zen also uses it as the PDF alt text of the first picture in that figure. Write it from the data: the chart type, the axes and their ranges, and the values or trends a reader needs (every figure in this manual has one; look at the source of \autoref{fig:bargraph}). If a figure has a picture but no description, the log shows a \texttt{CHIZen Warning}.
    \item \textbf{Mark the header rows of tables.} \texttt{acmart} assumes the first row of a \texttt{tabular} holds the column headers. If your table has two header rows, say so before it with \texttt{\textbackslash tagpdfsetup\{table/header-rows=\{1,2\}\}}. Keep every row the same width: add empty cells (\texttt{\&}) rather than ending a row early.
    \item \textbf{Write meaningful link text}, such as ``the \href{https://sigchi.org/resources/guides-for-authors/accessibility/}{SIGCHI Accessibility Guide}'' rather than ``click here''.
    \item \textbf{Get it right in black and white}: pair colour with labels, shapes or patterns, and check contrast.
    \item \textbf{Avoid subfigures} (\texttt{subcaption}) if you need a fully valid PDF/UA file: the tagging project lists \texttt{subcaption} as not yet compatible, and figures with subfigures fail PDF/UA validation (two such figures in this manual do). Packages such as \texttt{minted} (code) and \texttt{changes} (tracked changes) compile but are not yet tagged meaningfully.
    \item \textbf{Check the PDF.} The free \href{https://verapdf.org}{veraPDF} validator checks PDF/UA and PDF/A, and this repository's GitHub workflow runs it on every push. A validator cannot judge whether a description is useful: read your descriptions aloud, or listen to the paper with a screen reader.
\end{enumerate}


% %%%%%%%%%%%%%%%%%%%%%%%%%%%%%%%%%%%%%%%%%%%%%%%%%%%%%%%%%%%%%%%%%%%%%%%%%%%%%
% \section{\LaTeX\ Packages and Tools}
% %%%%%%%%%%%%%%%%%%%%%%%%%%%%%%%%%%%%%%%%%%%%%%%%%%%%%%%%%%%%%%%%%%%%%%%%%%%%%

% - Lots of different options here
% - Many come down to preference

% - have included quite a few here, as needs evolved over years
% - some basic documentation

\newpage
\begin{markdown}

# \LaTeX\ Packages and Tools

This template pre-loads a number of different packages by default. 

1. 

\end{markdown}

\newpage

\section{Code}
\label{sec:code}

If you need to share some code, I recommend using the \texttt{minted} package, which is already imported by CHI Zen. The \texttt{minted} package provides a nice environment for sharing pieces of code directly in your document, and has built-in support for a long list of common programming languages. It can highlight keywords, add line numbers, and comes with a variety of colour styles to fit your needs. It can also import code directly from a file, so you can link directly to your own source code if appropriate. 

If you don't need an entire code block, the package also enables you to insert single lines of code into your manuscript, like this: 

\mint{python}|fibonacci(100) # Defined in the figure below|





\begin{figure}[b!]
\Description{Two code listings side by side. Left, in a colour style similar to Visual Studio: a Python function, fibonacci(nterms), that starts a list with the terms 1 and 1, appends the sum of the previous two terms until it has nterms terms, and prints each term. Right, in a black-and-white style: a C main function that prints hello, world and returns 0. Both listings have line numbers and a rule above and below.}
    \centering
    \begin{subfigure}{.45\textwidth}
    \usemintedstyle{vs}
    \begin{minted}[
        frame=lines,
        framesep=2mm,
        baselinestretch=1.2,
        fontsize=\footnotesize,
        linenos]
        {python}
    # An Example in Python

    def fibonacci (nterms):

    terms = [1,1]

    for i in range(2,nterms):
        terms += [terms[i-1] + terms[i-2]]

    for term in terms:
        print(term)
    \end{minted}
    \caption{Example code in Python}
    \end{subfigure}
    %
    \begin{subfigure}{.45\textwidth}
    \usemintedstyle{bw}
    \begin{minted}[
        frame=lines,
        framesep=2mm,
        baselinestretch=1.2,
        fontsize=\footnotesize,
        linenos]
        {c}
// An Example in C
int main() {
    printf("hello, world");
return 0;
}
\end{minted}
\caption{Example code in C}
\end{subfigure}
    \caption{Example code fragments in \texttt{minted} for Python and C. The colour style on the left is similar to what you might see in an IDE, the one on the right is black and white, and might be better when printed. }
    \label{fig:minted}
\end{figure}




\subsection*{Additional Resources}

\begin{enumerate}
    \item The \texttt{minted} package on CTAN: \href{https://www.ctan.org/pkg/minted}{https://www.ctan.org/pkg/minted}
    \item Overleaf's \texttt{minted} documentation: \href{https://www.overleaf.com/learn/latex/Code\_Highlighting\_with\_minted}{https://www.overleaf.com/learn/latex/Code\_Highlighting\_with\_minted}
\end{enumerate}

%
% This Section is written in LaTeX because Markdown wasn't playing well with the tables
%
\newpage
%%%%%%%%%%%%%%%%%%%%%%%%%%%%%%%%%%%%%%%%%%%%%%%%%%%%%%
\section{Tables}
\label{sec:tables}
%%%%%%%%%%%%%%%%%%%%%%%%%%%%%%%%%%%%%%%%%%%%%%%%%%%%%%

As a rule of thumb, if you can show your data with a table, consider doing that first. Tables are quite effective at `showing your data', and are usually easy to interpret. The Vis literature even suggests that tables might be more convincing for a skeptical audience \citep{pandey2014persuasive}!

In terms of table formatting, many of the same principles apply to tables as figures. My advice is to: 

\begin{enumerate}
    \item Maintain a high Data-Ink Ratio \citep{tufte1983visual}: 
        \begin{itemize}
            \item Don't use vertical lines
            \item Use horizontal lines sparingly
            \item Use whitespace to help group things together that belong together
        \end{itemize}
    \item Use \texttt{\textbackslash toprule}, \texttt{\textbackslash midrule}, and \texttt{\textbackslash bottomrule} to frame your table
    \item Use the \texttt{[tb!]} flags when declaring the table in \LaTeX to push it to the top or bottom of a page
    \item Use the \texttt{label\{\}} command to give your table a label that you can use later with \texttt{autoref\{\}}
\end{enumerate}
%
In general, writing tables directly in \LaTeX is a huge pain --- the output is generally \emph{beautiful}, but the input is complicated, and generally not very human-readable. I suggest using an online table generator, like \href{https://www.tablesgenerator.com}{https://www.tablesgenerator.com}, and then fine-tuning the formatting in your document. \autoref{tab:exampletable} is an example from \citet{wallace2013}. 

\begin{center}
\begin{table*}[b!]
\small
\begin{center}
\begin{tabular}{l l l l l}
\toprule
Condition & \multicolumn{4}{l}{Teamwork Measures}\\
\cmidrule(r){2-5}
& \multicolumn{2}{l}{Measure 1} & Measure 2 & Measure 3 \\
\cmidrule(r){2-3} 
&  Sub 1 & Sub 2 &  \\ 


\midrule


Condition 1 & .478 & .468 & 1.68 & 12.3  \\
                  & (.0259) & (.031) & (.703) & (5.93) \\
\\

Condition 2 & .444 & .451 & 1.32 & 18.4 \\
                              & (.0478) & (.0585) & (.319) & (2.76) \\
\\

Condition 3 &  & .410 &  2.857 &  \\
                     &        & (.0384)  & (.350) & \\
\\


ANOVA Results & $F_{(1,12)} = 2.27,$ & $F_{(2,18)} = 2.76,$  & $F_{(2,18)} = 16.16,$ & $F_{(1,12)} = 6.16,$ \\
                          & $p = .158$ & $p = .090  $ &  $p < .0001$* &  $p = 0.029$ * \\
\bottomrule
\end{tabular}

\caption[a]{Mean values and standard deviations (in parentheses) for teamwork measures, and ANOVA results for comparisons between experimental conditions. Significant results denoted by *.}
\label{tab:exampletable}
\end{center}
\end{table*}
\end{center}

\subsection*{Additional Resources}
\begin{enumerate}
    \item Overleaf's documentation for tables: \href{https://www.overleaf.com/learn/latex/Tables}{https://www.overleaf.com/learn/latex/Tables}
\end{enumerate}


%
% ADAPT 
%


% \definecolor{palered}{HTML}{ffadad}
% \definecolor{paleorange}{HTML}{ffd6a5}
% \definecolor{paleyellow}{HTML}{fdffb6}
% \definecolor{palegreen}{HTML}{caffbf}
% \definecolor{palecyan}{HTML}{9bf6ff}
% \definecolor{paleblue}{HTML}{a0c4ff}
% \definecolor{palepurple}{HTML}{bdb2ff}
% \definecolor{palepink}{HTML}{ffc6ff}
% \definecolor{palewhite}{HTML}{fffffc}

% \begin{table}[bt!]
%     \centering
%     \small
    
% \begin{tabular}{p{4.5cm}p{3cm}p{3cm}p{3cm}}
% \toprule
% \textbf{Main Themes} & \textbf{Sub Themes} \\ 
% \midrule
% \rowcolor{palered!10}\textcolor{palered}{\faCircle}~\textbf{(1)} Decisions to play are based on childhood experiences relative to the cultural norm growing up  &  \textcolor{palered}{\faChevronCircleRight}~\emph{Nostalgic memories  and tabletop games}  & \textcolor{palered}{\faChevronCircleRight}~\emph{Families with low SES} & \\ 

% \rowcolor{paleorange!10}\textcolor{paleorange}{\faCircle}~\textbf{(2)} Adapted and influenced cultural themes in tabletop games should teach more about culture &  \textcolor{paleorange}{\faChevronCircleRight}~\emph{Participants call on game designers} & & \\ 

% \rowcolor{paleblue!10}\textcolor{paleblue}{\faCircle}~\textbf{(3)} Having someone in mind to play with is the strongest contribution to purchase decision of a tabletop game &  \textcolor{paleblue}{\faChevronCircleRight}~\emph{Interested players and limited players}  & \textcolor{paleblue}{\faChevronCircleRight}~\emph{Word of mouth through price, reviews, and demonstration} & \textcolor{paleblue}{\faChevronCircleRight}~\emph{Quality of game components} \\ 

% \rowcolor{palegreen!10}\textcolor{palegreen}{\faCircle}~\textbf{(4)} Bonding experiences emerge through collaborative and competitive playstyle regardless of household income &  \textcolor{palegreen}{\faChevronCircleRight}~\emph{Disengaged from the digital world} & & \\ 

% \rowcolor{palecyan!10}\textcolor{palecyan}{\faCircle}~\textbf{(5)} Portability is connected for playing  frequency of a tabletop game &  \textcolor{palecyan}{\faChevronCircleRight}~\emph{Playing space and size matters} & \textcolor{palecyan}{\faChevronCircleRight}~\emph{Storage space} &  \textcolor{palecyan}{\faChevronCircleRight}~\emph{High-SES families prefer games with a higher learning curve} \\ 
% \rowcolor{palepurple!10}\textcolor{palepurple}{\faCircle}~\textbf{(6)} Tabletop games should have app-assisted technology   & \textcolor{palepurple}{\faChevronCircleRight}~\emph{Tutorial that demonstrates how the game should be played} &  \textcolor{palepurple}{\faChevronCircleRight}~\emph{Game rules are confusing} &  \\ 

% \rowcolor{palepink!10}\textcolor{palepink}{\faCircle}~\textbf{(7)} Involvement with tabletop games community culture  &  \textcolor{palepink}{\faChevronCircleRight}~\emph{Community involvement} & & \\ 
% \bottomrule
% \end{tabular}

% \caption{We inferred seven themes from our hybrid (largely reflexive) thematic analysis (drawing from codebook approaches) that are relevant to cultural influences, play and purchase decisions in tabletop games. We present a complete overview of the themes and extrapolated subthemes in this table.}
% \label{tab:findings}
% \end{table}


\newpage

%
% This section written in LaTeX because a lot of the graphing is untested in Markdown
%


%%%%%%%%%%%%%%%%%%%%%%%%%%%%%%%%%%%%%%%%%%%%%%%%%%%%%%%%%%%%%%%%%%%%%%%%%%%%%
\section{Figures}
\label{sec:figures}
%%%%%%%%%%%%%%%%%%%%%%%%%%%%%%%%%%%%%%%%%%%%%%%%%%%%%%%%%%%%%%%%%%%%%%%%%%%%%

\begin{markdown}
Figures and Tables provide your readers with a `visual abstract' of your work. They can make the (often dense) results section more accessible to readers who are unfamiliar with your research area, and are a fallback when something in the text doesn't make sense right away. That is, they're often the first impression a reader has of your paper, make your work more accessible to the general audience, and provide a second, more visual, way of explaining results to your readers.

A general principal that I try to embody in all of these graphs is to **show your data**. Too often, we provide only summary data like an average and a standard error value, when we could have shown the entire data set. Summary statistics are useful, but often hide important trends in the data. (We should take this to heart well before writing the paper up, this all holds true for *data analysis*). I think that in the past, tools weren't always up to the task of showing our data. Tools like Excel would help you easily produce a bar chart, but don't have a preset for box plots. Fortunately, today our tools have become much better. 

I've also tried to take best practice into consideration where possible, but there's always room for improvement. \citet{munznervisualization} provides an excellent summary of best practice from the literature, and I encourage you to read Chapter 6 on Rules of Thumb in depth to start thinking about how to approach thinking about visualizations as you're working on them. In particular, tips like **get it right in black and white** and **eyes beat memory** will help to improve the readability *and* accessibility of your paper. 

The following sections are framed in terms of problem/solution for common types of graphs that we see in HCI research, with examples for each. The style for each graph is informed and inspired by \citet{tufte1983visual}, who emphasizes minimalist presentation of quantitative data. I find that Tufte's principles of information visualization, and general approach to working on graphs, are largely consistent with Munzner's advice above. That said, the style cetainly isn't for everyone --- your mileage may vary, please use what you find useful.  


 \end{markdown}
 
%  For example, some good examples:
% Bartneck and Hu \cite{Bartneck:2009:SAC:1518701.1518810}


% Senellart \cite{Senellart:2013:DYR:2513166.2514938} discusses some useful cases of `bad' style, which superficially may appear to be picky but offer important guidance towards submitting polished papers. 
 
 
% Munzer \cite{Munzner:2008:PPW:1422919.1422927} 
% - some good advice in general, but more focused on info vis research

% - text is large and legible
% - all axes are labelled, including units of measure
% - axes start at 0, unless explicitly justified (for a good reason!)


%%%%%%%%%%%%%%%%%%%%%%%%%%%%%%%%%%%%%%%%%%%%%%%%%%%%%%%%%%
\subsection{Colour Palettes}
\label{sec:palettes}
%%%%%%%%%%%%%%%%%%%%%%%%%%%%%%%%%%%%%%%%%%%%%%%%%%%%%%%%%%

\texttt{CHIZen/figures.tex} defines several palettes as PGFPlots cycle lists (\texttt{cycle list name=...}): \texttt{okabe-ito} (the default for the bar, stacked bar, panel and box plot styles), \texttt{IDEA lines} (the Okabe--Ito colours with dash patterns and marks; the default for \texttt{IDEA line}), \texttt{waterloo-light}, \texttt{waterloo-dark}, \texttt{colorbrewer-RYB}, and \texttt{colorbrewer-RYB-plain}, plus the divergent Likert colours. To use another palette, set \texttt{cycle list name} in your figure's axis options. To \emph{get it right in black and white}, keep in mind that:

\begin{itemize}
    \item \textbf{None of the categorical palettes is greyscale-safe.} In \texttt{waterloo-light}, ArtsLight and SchoolRedLight, and MathLight and ScienceLight, print as almost the same grey; most of \texttt{waterloo-dark} prints as one mid grey; and \texttt{colorbrewer-RYB} has pairs of near-identical lightness too. In greyscale, tell series apart with marks, dash patterns, or direct labels as well as colour (\texttt{IDEA lines} varies all three).
    \item \textbf{RYB2 (pale yellow) is almost invisible on white.} Use \texttt{colorbrewer-RYB}, which draws a darker outline around each fill, rather than \texttt{colorbrewer-RYB-plain}.
    \item \textbf{The default, \texttt{okabe-ito}, is colour-blind safe}: the Okabe--Ito palette stays distinguishable with the common forms of colour blindness. It is not greyscale-safe either (its orange and sky blue print as the same grey), so the same advice about marks and labels applies. Its yellow is light on white, so it is drawn with a darker outline.
\end{itemize}







\newpage
%%%%%%%%%%%%%%%%%%%%%%%%%%%%%%%%%%%%%%%%%%%%%%%%%%%%%%%%%%
\subsection{Comparing Performance: Bar Charts}
%%%%%%%%%%%%%%%%%%%%%%%%%%%%%%%%%%%%%%%%%%%%%%%%%%%%%%%%%%

Bar graphs are probably the most commonly-used data visualization at CHI. They are simple, easy to create, and easy to understand. And so they make comparisons between two or more variables easy, since data is presented in a side=by-side manner on a common scale (the y-axis). Because of these advantages, they are often used to make quantitative comparisons based on criteria like performance, efficiency, and error rate. A very common approach to presenting data in a bar chart is to show means and standard error  (e.g., \autoref{fig:bargraph} from \citet{pietroszek2015}).






\begin{figure}[b!]
\Description{Bar chart of mean selection time in seconds (y axis, 0 to 10) by occlusion (x axis: non-occluded and occluded targets), with an error bar on each bar. Non-occluded targets: Smartcasting 2.10 s and Tiltcasting 3.24 s. Occluded targets: Tiltcasting 4.0 s and Depth Cursor 6.8 s. Depth Cursor has the longest bar and the largest error bar.}
\begin{tikzpicture}
\begin{axis}[IDEA bar,
        symbolic x coords={non-occluded,occluded},
        xtick = {non-occluded,occluded},
        enlarge x limits=0.5,
        xlabel = {Occlusion},
	    ymin = 0, ymax = 10,
        ylabel = {Selection Time (s)},
        width=\columnwidth,
        height = 5cm,
]

%Depth Ray
\addplot[style={draw=none,fill=AHSLight},error bars/.cd, y dir = both, y explicit, error bar style={black}]
coordinates {(occluded, 6.8)  +- (1.7,1.7)};
\addlegendentry{Depth Cursor}

%Tiltcasting
\addplot+[style={draw=none, fill=SchoolRedLight,line width = 0pt},error bars/.cd, y dir = both, y explicit, error bar style={black}]
coordinates {(non-occluded,3.24) +- (.712,.712)
		    (occluded,4)  +- (.922,.922)};
\addlegendentry{Tiltcasting}

%Smartcasting
\addplot+[style={draw=none, fill=EnvironmentLight},error bars/.cd, y dir = both, y explicit, error bar style={black}]
coordinates {(non-occluded, 2.10)  +- (.42,.42)};
\addlegendentry{Smartcasting}

\end{axis}
\end{tikzpicture}
\caption{An example bar graph with error bars. Data is specified directly in the latex file --- not in a separate data file.}
\label{fig:bargraph}
\end{figure}




\begin{figure*}[tb!]
\Description{Bar chart of agreement rate (y axis, 0 to 1) for each of 34 gesture tasks (x axis, labels rotated), sorted from highest to lowest agreement. Bar shade shows the agreement category: three tasks have high agreement (Play Song 0.43, Jump to Movie 0.34, Jump to Floor Level 0.30), 16 tasks from Play Movie (0.26) to Jump to Store (0.11) have medium agreement, and 15 tasks from Toggle View (0.08) to Jump to Menu Item (0.01) have low agreement. An orange segment stacked on top of each bar shows the agreement rate when gestures are grouped by type and length; it ranges from 0.03 (Play Song) to 0.61 (Jump to Flight) and is larger than the visual-trace agreement for 29 of the 34 tasks.}
%\centerline{
%\includegraphics[width=1.75\columnwidth]{figures/gesture-agreement.png}}
\begin{tikzpicture}
\begin{small}
\begin{axis}[IDEA stacked bar, 
	bar width = 5pt,
	% X Axis options (xmin, xmax, xtick={})
    xtick={1,2,3,4,5,6,7,8,9,10,11,12,13,14,15,16,17,18,19,20,21,22,23,24,25,26,27,28,29,30,31,32,33,34},
    xticklabels={Play Song, Jump to Movie, Jump to Floor Level, Play Movie, Pause Song, Replay Song, Activate Search, Zoom In, Zoom Out, Favorite Item, Activate Search, Next Movie, Previous Movie, Pan Up, Pan Down, Pan Left, Pan Right, Jump to Flight, Jump to Store, Toggle View, Home Screen, Floor Up, Floor Down, Next Song, Previous Song, Page Left, Page Right, Reset Location, Clear Filter, Show Menu, Volume up, Volume Down, Filter Delayed Flights, Jump to Menu Item},
    xlabel={Task},
    %xlabel near ticks, 
    xticklabel style={rotate=45,anchor=north east,font=\small},
	% Y Axis options (ymin, ymax, ytick={})
    ymin = 0, ymax = 1.0,
    ylabel={Agreement Rate},
    ytick={0,.25,.5,.75,1.0},
    %general sizing options
    width = \columnwidth,
    height = 4cm,
	enlarge x limits = {abs=0.25cm},
    % legend
    legend image code/.code={\draw[#1] rectangle (0.25cm,0.25cm);},
    legend style={at={(xlabel.south)},yshift=-0.5ex, xshift=-2ex, anchor=north, legend columns = 4, draw=none, column sep = 0.25cm},
       ]



    % DATA GOES IN HERE
    

    \addplot+[style={draw=ScienceDark,fill=ScienceDark}] plot coordinates {
    (1,0.43)
	(2,0.34)
	(3,0.30)
    (4, 0)
	(5, 0)
	(6, 0)
	(7, 0)
	(8, 0)
	(9, 0)
	(10, 0)
	(11, 0)
	(12,0)
	(13,0)
	(14,0)
	(15,0)
	(16,0)
	(17,0)
	(18,0)
	(19,0)
    (20, 0)
	(21, 0)
	(22, 0)
	(23, 0)
	(24, 0)
	(25, 0)
	(26, 0)
	(27, 0)
	(28, 0)
	(29, 0)
	(30, 0)
	(31, 0)
	(32, 0)
	(33, 0)
	(34, 0)
};
\addlegendentry{High ($0.3 - 0.5$)}
    
    \addplot+[style={draw=ScienceDark!75!white,fill=ScienceDark!75!white}] plot coordinates 
{
    (1,0)
	(2,0)
	(3,0)
    (4, 0.26)
	(5, 0.26)
	(6, 0.23)
	(7, 0.23)
	(8, 0.21)
	(9, 0.21)
	(10, 0.20)
	(11, 0.19)
	(12,0.17)
	(13,0.17)
	(14,0.15)
	(15,0.15)
	(16,0.15)
	(17,0.15)
	(18,0.12)
	(19,0.11)
    (20, 0)
	(21, 0)
	(22, 0)
	(23, 0)
	(24, 0)
	(25, 0)
	(26, 0)
	(27, 0)
	(28, 0)
	(29, 0)
	(30, 0)
	(31, 0)
	(32, 0)
	(33, 0)
	(34, 0)
};
\addlegendentry{Medium ($0.1 - 0.3$)}
    
    \addplot+[style={draw=ScienceLight,fill=ScienceLight}] plot coordinates 
{
(1,0)
	(2,0)
	(3,0)
    (4, 0)
	(5, 0)
	(6, 0)
	(7, 0)
	(8, 0)
	(9, 0)
	(10, 0)
	(11, 0)
	(12,0)
	(13,0)
	(14,0)
	(15,0)
	(16,0)
	(17,0)
	(18,0)
	(19,0)
    (20, 0.08)
	(21, 0.06)
	(22, 0.06)
	(23, 0.06)
	(24, 0.06)
	(25, 0.06)
	(26, 0.05)
	(27, 0.05)
	(28, 0.04)
	(29, 0.03)
	(30, 0.03)
	(31, 0.03)
	(32, 0.03)
	(33, 0.02)
	(34, 0.01)
};
\addlegendentry{Low ($\leq 0.1$)}
\addplot+[style={draw=ArtsLight!60!white,fill=ArtsLight!60!white}] plot coordinates {
    (1,0.03)
	(2,0.04)
	(3,0.03)
	(4,0.34)
	(5,0.27)
	(6,0.16)
	(7,0.17)
	(8,0.59)
	(9,0.59)
	(10,0.37)
	(11,0.21)
	(12,0.55)
	(13,0.55)
	(14,0.35)
	(15,0.35)
	(16,0.35)
	(17,0.35)
	(18,0.61)
	(19,0.38)
    (20, 0.23)
	(21, 0.37)
	(22, 0.31)
	(23, 0.31)
	(24, 0.35)
	(25, 0.35)
	(26, 0.44)
	(27, 0.44)
	(28, 0.25)
	(29, 0.40)
	(30, 0.19)
	(31, 0.34)
	(32, 0.34)
	(33, 0.16)
	(34, 0.20)
};
\addlegendentry{Agreement by Type and Length}
\end{axis}
\end{small}
\end{tikzpicture}

\caption{Agreement Rates (AR) for suggested gestures in terms of its visual trace, by each of the 34 tasks included in our elicitation study. Gestures ranked by agreement rating, with Vatavu and Wobbrock's high, medium, and low  agreement categories indicated by color. Orange bars show agreement scores when gestures are placed within more general groups of type and length.}
\label{fig:gesture-agreement}
\end{figure*}





\newpage
%%%%%%%%%%%%%%%%%%%%%%%%%%%%%%%%%%%%%%%%%%%%%%%%%%%%%%%%%%
\subsection{Comparing Performance: Box Plots, Violin Plots, Bee Swarm Plots}
%%%%%%%%%%%%%%%%%%%%%%%%%%%%%%%%%%%%%%%%%%%%%%%%%%%%%%%%%%


- to be more accurate and descriptive, box plots are often used
- similar to a bar graph with error bars, but also indicate quartiles and outliers 
- typically, better practice to include a box plot when possible -- but can be substantially more difficult to create in Microsoft Word. 

variations: 
- Admittedly, "Bee Swarm" plots are my personal preference for this kind of data since they "show your data" in a very clear way... but I haven't yet found a  way to generate them in \LaTeX. If you happen to work in R, there are some excellent supports for generating these graphs, and you can import those directly into this template with a little work. 
- Violin plots 


%%%%%%%%%%%%%%%%%%%%%%%%%%%%%%%%%%%%%%%%%%%%%%%%%%%%%%%%%%%%%%%%%%
% Gigantic Figure Starts here for physiological data... 
%%%%%%%%%%%%%%%%%%%%%%%%%%%%%%%%%%%%%%%%%%%%%%%%%%%%%%%%%%%%%%%%%%

\definecolor{Nature}{HTML}{3AA876}
\definecolor{Urban}{HTML}{D14820}


\pgfplotstableread{
Treatment    P0_HR    P0_GSR    P1_HR    P1_GSR    P2_HR    P2_GSR
0	-0.06	-1.09	0.73	-0.54	-0.57	0.87
0	0.38	-1.022	0.63	0.32	-0.88	0.106666667
0	-0.52	-1.31	0.59	0.47	-0.72	-0.36
0	-0.49	-1.37	0.99	0.01	-0.27	0.32
0	0.2	-1.022	-0.2	0.32	0.11	0.106666667
0	0.76	1.14	0.15	0.46	-0.45	-0.91
0	0.3	-1.022	0.33	0.32	-0.67	0.106666667
0	-0.73	-0.44	1.13	0.41	-0.54	-0.12
0	-0.62	-1.04	0.94	0.74	-0.69	-0.33
0	-0.66	-2.46	0.63	0.23	-0.68	0.52
}\SPN 

\pgfplotstableread{
Treatment    P0_HR    P0_GSR    P1_HR    P1_GSR    P2_HR    P2_GSR
1	-0.38	-1.022	0.51	0.32	-0.66	0.106666667
1	-0.45	-0.98	0.44	0.81	-0.55	-0.97
1	0.31	-0.8	0.78	0.74	-0.81	0.31
1	-0.35	-0.89	0.75	0.87	-0.56	-0.44
1	0.08	-0.8	0.7	1.46	-0.44	-0.47
1	-0.22	-1.7	0.64	-0.02	-0.71	0.69
1	-0.29	-0.77	0.61	-0.52	-0.74	0.84
1	-1.24	-2.24	0.42	-0.06	-0.19	0.63
1	-0.06	-1.022	0.25	0.32	-0.25	0.106666667
1	0.14	-0.58	0.57	-0.26	-0.86	1.02
}\SPU

\pgfplotstableread{
Treatment    Z0_posaffect    Z0_negaffect    Z0_attentive    Z1_posaffect    Z1_negaffect    Z1_attentive    Z2_posaffect    Z2_negaffect    Z2_attentive
0	3.75	1.22	3.50	2.00	3.56	2.50	3.75	1.22	3.50
0	1.50	1.11	3.00	1.00	1.44	2.50	2.25	1.00	2.50
0	3.25	1.00	3.50	1.75	2.11	1.50	3.00	1.11	4.00
0	2.50	1.00	3.50	3.00	1.11	3.50	3.00	1.22	4.00
0	2.00	1.00	2.50	1.00	1.11	3.00	2.00	1.11	2.00
0	2.50	1.22	4.00	1.00	2.22	1.50	2.25	1.11	3.50
0	3.50	1.33	2.50	2.50	2.33	1.00	2.25	2.67	1.50
0	1.00	1.11	2.00	1.25	1.33	2.00	1.50	1.00	1.50
0	3.25	1.11	2.50	2.00	1.67	3.50	4.25	1.00	3.50
0	3.00	1.00	3.00	1.75	1.78	2.00	2.50	1.11	2.50
}\SZN

\pgfplotstableread{
Treatment    Z0_posaffect    Z0_negaffect    Z0_attentive    Z1_posaffect    Z1_negaffect    Z1_attentive    Z2_posaffect    Z2_negaffect    Z2_attentive
1	4.00	1.00	3.00	3.25	1.44	2.00	4.75	1.00	2.50
1	2.75	1.22	2.50	1.00	4.22	1.50	2.25	1.33	3.00
1	2.75	1.11	2.50	1.00	1.72	2.00	1.00	1.78	1.00
1	2.25	1.56	3.50	1.50	2.33	3.00	3.25	1.00	3.50
1	2.00	1.11	3.50	1.25	3.00	3.00	2.50	1.11	3.00
1	3.75	1.33	4.00	1.25	2.78	3.50	3.25	1.22	3.00
1	3.25	1.00	3.50	2.25	1.44	3.50	2.75	1.11	1.00
1	3.50	1.11	3.50	1.25	1.67	4.00	3.75	1.22	2.50
1	3.25	1.44	4.00	3.00	2.33	3.50	3.75	1.33	3.50
1	2.25	1.00	2.50	1.50	3.22	2.00	3.75	1.00	1.50
}\SZU

\begin{figure}[b!]
\Description{Two box plots of z-scores at Baseline, Stressor and Game, each with Nature (green) and Urban (orange) boxes side by side. Heart rate: medians rise from about -0.3 at Baseline to about 0.6 at Stressor and fall to about -0.6 during the Game in both conditions. Galvanic skin response: medians rise from about -1.0 at Baseline to about 0.3 at Stressor and are about 0.1 (Nature) and 0.2 (Urban) during the Game. Nature and Urban boxes overlap at every stage.}
\begin{tikzpicture}
\begin{axis}[
    IDEA box,
    height = 6cm,
    width = \textwidth,
    ymin=-2.75, ymax=2.75,
    ylabel = Heart Rate ($z$-score), 
    ytick={-2,-1,0,1,2},
    xtick={1,2,3},
    xticklabels={Baseline,Stressor,Game},
]

%% Baseline
\addplot[Urban, mark=*,
	boxplot,
    forget plot,
    rshift,
]
table[y=P0_HR] {\SPU};
    

\addplot[Nature,mark=*,
    lshift,
    boxplot,
]
table[y=P0_HR] {\SPN};


%% Stressor
\addplot[Urban, mark=*,
	boxplot,
    forget plot,
    rshift,
]
table[y=P1_HR] {\SPU};
    

\addplot[Nature,mark=*,
    lshift,
    boxplot,
]
table[y=P1_HR] {\SPN};


%% GAME
\addplot[Urban, mark=*,
	boxplot,
    forget plot,
    rshift,
]
table[y=P2_HR] {\SPU};

    
\addplot[Nature,mark=*,
    lshift,
    boxplot,
]
table[y=P2_HR] {\SPN};

   
\end{axis}
\end{tikzpicture}
%\caption{Box plot of heart rate data collected during the study. Participants in Nature (left) vs Urban (right) conditions.}
%\end{figure}


%\begin{figure}
\begin{tikzpicture}
\begin{axis}[
    IDEA box,
	height= 6cm,
    width = \textwidth,
    ymin=-2.75, ymax=2.75,
    ylabel = Galvanic Skin Response ($z$-score), 
    ytick={-2,-1,0,1,2},
    xtick={1,2,3},
    xticklabels={Baseline,Stressor,Game},
    legend style={at={(0.25,-.2)},anchor=south,font={\sffamily}}
]

%% Baseline
\addplot[Urban, mark=*,
	boxplot,
    forget plot,
    rshift,
]
table[y=P0_GSR] {\SPU};
    

\addplot[Nature,mark=*,
    lshift,
    boxplot,
]
table[y=P0_GSR] {\SPN};


%% Stressor
\addplot[Urban, mark=*,
	boxplot,
    forget plot,
    rshift,
]
table[y=P1_GSR] {\SPU};
    

\addplot[Nature,mark=*,
    lshift,
    boxplot,
]
table[y=P1_GSR] {\SPN};
%\addlegendimage{fill=Nature,draw=Nature}
%\addlegendentry{Nature}

%% GAME
\addplot[Nature,mark=*,
    lshift,
    forget plot,
    boxplot,
]
table[y=P2_GSR] {\SPN};
%\fill[Nature] (0,0) circle (0.1cm);


\addplot[Urban, mark=*,
	boxplot,
    rshift,
]
table[y=P2_GSR] {\SPU};
%\fill[Urban] (0,0) circle (0.1cm);
%\addlegendimage{fill=Urban,draw=Urban}
%\addlegendentry{Urban}
    
    \coordinate (legend) at (axis description cs:.15,-.15);
\end{axis}
% this is a dummy `axis' environment only to create the legend
\matrix [
    matrix of nodes, 
    every node/.style={anchor=center}, 
    ] at (legend) {
        |[fill=Nature]| & |[font=\small\sffamily]|Nature & |[fill=white]| & |[fill=Urban]| & |[font=\small\sffamily]|Urban \\
    };
\end{tikzpicture}
\caption{A summary of Heart Rate and Galvanic Skin Response data collected during the study. Our analysis revealed a significant increase between Baseline and Stressor data, and a significant decrease between Stressor and Game data for both measures, as anticipated from our experimental design. We found no significant differences between the Nature and Urban game settings.}
\label{fig:physio}
\end{figure}

\begin{figure}[b!]
\Description{Three box plots of self-reported ratings on a 1 to 5 scale at Baseline, Stressor and Game, each with Nature (green) and Urban (orange) boxes side by side. Positive affect: medians drop at Stressor (Nature from about 2.8 to 1.8, Urban from 3.0 to 1.4) and recover during the Game (about 2.4 and 3.3). Negative affect: medians rise at Stressor (Nature from about 1.1 to 1.7, Urban from 1.1 to 2.3) and return to about 1.1 during the Game. Attentiveness: medians are about 3.0, 2.3 and 3.0 for Nature and 3.5, 3.0 and 2.8 for Urban. Nature and Urban boxes overlap at every stage.}
\begin{tikzpicture}
\begin{axis}[
    IDEA box,
    height=4cm,
    width = \textwidth,
    ymin=1, ymax=5,
    ylabel = Positive Affect, 
    ytick={1,2,3,4,5},
    xtick={1,2,3},
    xticklabels={Baseline,Stressor,Game},
]

%% Baseline
\addplot[Urban, mark=*,
	boxplot,
    forget plot,
    rshift,
]
table[y=Z0_posaffect] {\SZU};
    

\addplot[Nature,mark=*,
    lshift,
    boxplot,
]
table[y=Z0_posaffect] {\SZN};


%% Stressor
\addplot[Urban, mark=*,
	boxplot,
    forget plot,
    rshift,
]
table[y=Z1_posaffect] {\SZU};
    

\addplot[Nature,mark=*,
    lshift,
    boxplot,
]
table[y=Z1_posaffect] {\SZN};


%% GAME
\addplot[Urban, mark=*,
	boxplot,
    forget plot,
    rshift,
]
table[y=Z2_posaffect] {\SZU};
    

\addplot[Nature,mark=*,
    lshift,
    boxplot,
]
table[y=Z2_posaffect] {\SZN};
   
\end{axis}
\end{tikzpicture}
\begin{tikzpicture}
\begin{axis}[
    IDEA box,
    height=4cm,
    width = \textwidth,
    ymin=1, ymax=5,
    ylabel = Negative Affect, 
    ytick={1,2,3,4,5},
    xtick={1,2,3},
    xticklabels={Baseline,Stressor,Game},
]

%% Baseline
\addplot[Urban, mark=*,
	boxplot,
    forget plot,
    rshift,
]
table[y=Z0_negaffect] {\SZU};
    

\addplot[Nature,mark=*,
    lshift,
    boxplot,
]
table[y=Z0_negaffect] {\SZN};


%% Stressor
\addplot[Urban, mark=*,
	boxplot,
    forget plot,
    rshift,
]
table[y=Z1_negaffect] {\SZU};
    

\addplot[Nature,mark=*,
    lshift,
    boxplot,
]
table[y=Z1_negaffect] {\SZN};


%% GAME
\addplot[Urban, mark=*,
	boxplot,
    forget plot,
    rshift,
]
table[y=Z2_negaffect] {\SZU};
    

\addplot[Nature,mark=*,
    lshift,
    boxplot,
]
table[y=Z2_negaffect] {\SZN};
   
\end{axis}
\end{tikzpicture}
\begin{tikzpicture}
\begin{axis}[
    IDEA box,
    height=4cm,
    width = \textwidth,
    ymin=1, ymax=5,
    ylabel = Attentiveness, 
    ytick={1,2,3,4,5},
    xtick={1,2,3},
    xticklabels={Baseline,Stressor,Game},
]

%% Baseline
\addplot[Urban, mark=*,
	boxplot,
    forget plot,
    rshift,
]
table[y=Z0_attentive] {\SZU};
    

\addplot[Nature,mark=*,
    lshift,
    boxplot,
]
table[y=Z0_attentive] {\SZN};


%% Stressor
\addplot[Urban, mark=*,
	boxplot,
    forget plot,
    rshift,
]
table[y=Z1_attentive] {\SZU};
    

\addplot[Nature,mark=*,
    lshift,
    boxplot,
]
table[y=Z1_attentive] {\SZN};


%% GAME
\addplot[Urban, mark=*,
	boxplot,
    forget plot,
    rshift,
]
table[y=Z2_attentive] {\SZU};
    

\addplot[Nature,mark=*,
    lshift,
    boxplot,
]
table[y=Z2_attentive] {\SZN};
   
\end{axis}
\end{tikzpicture}
\caption{A summary of Positive Affect, Negative Affect, and Attentiveness self-reported ratings collected during the study. All measures were aggregated across 5-point Likert scale questions from the ZIPERS questionnaire. As with the physiological measures (\autoref{fig:physio}), stress increased after the Stressor math test and decreased during gameplay, and we found no significant differences between the Nature and Urban game settings.}
\label{fig:selfreports}
\end{figure}

%%%%%%%%%%%%%%%%%%%%%%%%%%%%%%%%%%%%%%%%%%%%%%%%%%%%%%%%%%%%%%%%%%
% Gigantic Figure Ends here for physiological data... 
%%%%%%%%%%%%%%%%%%%%%%%%%%%%%%%%%%%%%%%%%%%%%%%%%%%%%%%%%%%%%%%%%%



\newpage
%%%%%%%%%%%%%%%%%%%%%%%%%%%%%%%%%%%%%%%%%%%%%%%%%%%%%%%%%%
\subsection{Likert Scale Data: Stacked Divergent Bar Graphs}
%%%%%%%%%%%%%%%%%%%%%%%%%%%%%%%%%%%%%%%%%%%%%%%%%%%%%%%%%%

A \emph{lot} of HCI research involves questionnaires, and interpreting diferences in questionnaire respones. For instance, Likert scale questions are very common. My preference for visualizing Likert scale responses is a stacked divergent bar graph (e.g., \autoref{fig:avatar_identification}), where responses are stacked into a bar, with neutral responses lined up at the graph's origin. This type of graph has a few properties that make it useful: in particular, readers can get a quick sense of `how positive' or `how negative' responses were based on the bar's position, while still being able to look at how many responses fell into each category. 

\begin{figure}[b!]
\Description{Diverging stacked bar chart of 7-point Likert responses to six avatar identification items, with two bars per item: Cleric (shades from blue for disagree to green for agree) and Monster (shades from red for disagree to purple for agree). Disagreement extends to the left of a centre line and agreement to the right, with neutral responses split across it; the x axis is the percentage of responses, 0 to 100 in each direction. Percentage agreeing, Cleric versus Monster: Connectedness 61 versus 28; Refer to avatar as we 61 versus 17; Physical similarity 22 versus 11; Personality similarity 78 versus 72; Physical ideal 44 versus 17; Ideal personality 83 versus 50. Asterisks mark Connectedness and Refer to avatar as we.}
\begin{tikzpicture}
\pgfplotstableread[col sep = comma]{data/Cleric_AvatarIdentificationStacked.csv}\ClericIdentificationData
\pgfplotstableread[col sep = comma]{data/Monster_AvatarIdentificationStacked.csv}\MonsterIdentificationData

\begin{axis}[
    IDEA Likert,
    height = 6.5cm,
    width = .75\columnwidth,
    % Y AXIS
    symbolic y coords = {Ava_Conn,	Ava_WE,	AvavsPhys,	AvaVsPers,	AvaVsPhysIdeal,	AvaVsPersIdeal},
    ytick={Ava_Conn, Ava_WE, AvavsPhys, AvaVsPers, AvaVsPhysIdeal, AvaVsPersIdeal},
    yticklabels={Connectedness, Refer to \\ Avatar as ``We'', Physical \\ Similarity, Personality\\ Similarity, Physical Ideal, Ideal Personality},
    yticklabel style={align=right,font={\sffamily\small}},
    y axis line style={draw=none},
    y dir = reverse,
    ymajorgrids = false,
    % X AXIS
    xmin=-100, xmax=100,
    xlabel={\% of Responses},
    xlabel style = {font=\small\sffamily},
    xticklabel style = {font=\small\sffamily},
    xtick = {-100,-50,0,50,100},
    xticklabels = {100, 50, 0, 50, 100},
    extra x ticks = {-75, 75},
    extra x tick labels = {$\leftarrow$ Disagree, Agree $\rightarrow$},
    every extra x tick/.style={major tick length=0,yshift={-8pt}},
    ]
    \addplot[draw=none,fill=none, forget plot] coordinates {(0,Ava_Conn)(0,Ava_WE)(0,AvavsPhys)(0,AvaVsPers)(0,AvaVsPhysIdeal)(0,AvaVsPersIdeal)};
    \addplot[draw=none,fill=Cleric_3, forget plot, bar shift=.1cm] table [x expr={\thisrow{4}*-0.5},y=Measure] {\ClericIdentificationData};
    \addplot[draw=none,fill=Cleric_2, forget plot, bar shift=.1cm] table [x expr={\thisrow{3}*-1},y=Measure] {\ClericIdentificationData};
    \addplot[draw=none,fill=Cleric_1, forget plot, bar shift=.1cm] table [x expr={\thisrow{2}*-1}, y=Measure ] {\ClericIdentificationData};
    \addplot[draw=none,fill=Cleric_0, forget plot, bar shift=.1cm] table [x expr={\thisrow{1}*-1}, y=Measure ] {\ClericIdentificationData};
    \resetstackedplots
    \addplot[draw=none,fill=none, forget plot] coordinates {(0,Ava_Conn)(0,Ava_WE)(0,AvavsPhys)(0,AvaVsPers)(0,AvaVsPhysIdeal)(0,AvaVsPersIdeal)};
    \addplot[draw=none,fill=Monster_3, forget plot, bar shift=-.1cm] table [x expr={\thisrow{4}*-0.5},y=Measure] {\MonsterIdentificationData};
    \addplot[draw=none,fill=Monster_2, forget plot, bar shift=-.1cm] table [x expr={\thisrow{3}*-1},y=Measure] {\MonsterIdentificationData};
    \addplot[draw=none,fill=Monster_1, forget plot, bar shift=-.1cm] table [x expr={\thisrow{2}*-1}, y=Measure ] {\MonsterIdentificationData};
    \addplot[draw=none,fill=Monster_0, forget plot, bar shift=-.1cm] table [x expr={\thisrow{1}*-1}, y=Measure ] {\MonsterIdentificationData};
\end{axis}

\begin{axis}[
    IDEA Likert,
    height = 6.5cm,
    width = .75\columnwidth,
    % Y AXIS
    symbolic y coords = {Ava_Conn,	Ava_WE,	AvavsPhys,	AvaVsPers,	AvaVsPhysIdeal,	AvaVsPersIdeal},
    ytick={Ava_Conn, Ava_WE, AvavsPhys, AvaVsPers, AvaVsPhysIdeal, AvaVsPersIdeal},
    yticklabels={Connectedness, Refer to \\ Avatar as ``We'', Physical \\ Similarity, Personality\\ Similarity, Physical Ideal, Ideal Personality},
    yticklabel style={align=right,font={\sffamily\small}},
    y axis line style={draw=none},
    y dir = reverse,
    ymajorgrids = false,
    % X AXIS
    xmin=-100, xmax=100,
    xlabel={\% of Responses},
    xlabel style = {font=\small\sffamily},
    xticklabel style = {font=\small\sffamily},
    xtick = {-100,-50,0,50,100},
    xticklabels = {100, 50, 0, 50, 100},
    extra x ticks = {-75, 75},
    extra x tick labels = {$\leftarrow$ Disagree, Agree $\rightarrow$ },
    every extra x tick/.style={major tick length=0,yshift={-8pt}},
    ]
    \addplot[draw=none,fill=none, forget plot] coordinates {(0,Ava_Conn)(0,Ava_WE)(0,AvavsPhys)(0,AvaVsPers)(0,AvaVsPhysIdeal)(0,AvaVsPersIdeal)};
    
    \addplot[draw=none,fill=Cleric_3, forget plot, bar shift=.1cm] table [x expr={\thisrow{4}*0.5}, y=Measure ] {\ClericIdentificationData};
    \addplot[draw=none,fill=Cleric_4, forget plot, bar shift=.1cm] table [x expr={\thisrow{5}}, y=Measure ] {\ClericIdentificationData};
    \addplot[draw=none,fill=Cleric_5, forget plot, bar shift=.1cm] table [x expr={\thisrow{6}}, y=Measure ] {\ClericIdentificationData};
    \addplot[draw=none,fill=Cleric_6, forget plot, bar shift=.1cm] table [x expr={\thisrow{7}}, y=Measure ] {\ClericIdentificationData}; \label{plot:cleric6}
    \resetstackedplots
    \addplot[draw=none,fill=none, forget plot] coordinates {(0,Ava_Conn)(0,Ava_WE)(0,AvavsPhys)(0,AvaVsPers)(0,AvaVsPhysIdeal)(0,AvaVsPersIdeal)};
    \addplot[draw=none,fill=Monster_3, forget plot, bar shift=-.1cm] table [x expr={\thisrow{4}*0.5}, y=Measure ] {\MonsterIdentificationData};
    \addplot[draw=none,fill=Monster_4, forget plot, bar shift=-.1cm] table [x expr={\thisrow{5}}, y=Measure ] {\MonsterIdentificationData};
    \addplot[draw=none,fill=Monster_5, forget plot, bar shift=-.1cm] table [x expr={\thisrow{6}}, y=Measure ] {\MonsterIdentificationData};
    \addplot[draw=none,fill=Monster_6, forget plot, bar shift=-.1cm] table [x expr={\thisrow{7}}, y=Measure ] {\MonsterIdentificationData};
    %\addlegendimage{only marks, mark=o}
    %\addlegendimage{only marks, mark=o}
    %\legend{Neutral, Agree, Strongly Agree, Disagree,Strongly Disagree}
    after end axis/.code={
        \node at (axis cs:80,Ava_Conn) [anchor=east, ,font=\tiny] {\textbf{*}};
        \node at (axis cs:80,Ava_WE) [anchor=east, font=\tiny] {\textbf{*}};
        %\draw [white] (axis cs:0,AvaVsPersIdeal) -- (axis cs:0,Ava_Conn);
        %\draw [white] (axis cs:-50,AvaVsPersIdeal) -- (axis cs:-50,Ava_Conn);
        %\draw [white] (axis cs:50,AvaVsPersIdeal) -- (axis cs:50,Ava_Conn);
    }
    \coordinate (legend) at (axis description cs:1.15,.85);
\end{axis}

% this is a dummy `axis' environment only to create the legend
\matrix [
    matrix of nodes, 
    every node/.style={anchor=center}, 
    ] at (legend) {
        |[fill=Cleric_0]| & |[fill=Cleric_1]| & |[fill=Cleric_2]| & |[fill=Cleric_3]| & |[fill=Cleric_4]| & |[fill=Cleric_5]| & |[fill=Cleric_6]| & |[font=\small\sffamily]|Cleric \\
        |[fill=Monster_0]| & |[fill=Monster_1]| & |[fill=Monster_2]| & |[fill=Monster_3]| & |[fill=Monster_4]| & |[fill=Monster_5]| & |[fill=Monster_6]| & |[font=\small\sffamily]|Monster \\
    };
\end{tikzpicture}
    
\caption{Summary of avatar identification responses, based on 7-point Likert scales. Mann Whitney U tests revealed significant differences between \textsc{Avatar} groups for `Connectness' and `Refer to Avatar as We'. }
\label{fig:avatar_identification}
\end{figure}





\newpage
%%%%%%%%%%%%%%%%%%%%%%%%%%%%%%%%%%%%%%%%%%%%%%%%%%%%%%%%%%
\subsection{Trends over Time: Line Graph}
%%%%%%%%%%%%%%%%%%%%%%%%%%%%%%%%%%%%%%%%%%%%%%%%%%%%%%%%%%



\begin{figure}[b!]
\Description{Two line charts of monthly counts from January 2014 to December 2017, with distinct users as a solid blue line and active threads as a dashed orange line. r/stopdrinking (left): distinct users rise from about 1,500 to about 6,000 and active threads from about 1,500 to about 5,400. r/OpiatesRecovery (right): distinct users rise from about 360 to about 770, peaking near 810, and active threads from about 280 to about 380, peaking near 480. Both lines in both charts trend upward with month-to-month fluctuation.}
\begin{tikzpicture}[baseline]
\begin{axis}[IDEA line,
	% Graph Options
	no markers, % dense monthly data: the dash pattern tells the series apart
	width=.475*\columnwidth,
	height=4cm,
	title=\texttt{r/stop\-drinking},
	% X Axis Options
	%date coordinates in=x,
	%date ZERO=2014-01-01,
    %xticklabel=\year,
    %xmin=2014-01-01,
    %xmax=2018-01-01,
    xlabel=Year,
	xtick={1,13,25,37,49},                              %% Jim: apologies... this is a hack
    xticklabels={2014,2015,2016,2017,2018},             %%      will figure out what's going on here if we have time later
	%xtick ={1,2,3,4,5,6,7,8,9},
	% Y Axis Options
	ylabel = Distinct Users / Active Threads,
	ylabel style = {font=\tiny},
	ymin = 0, ymax = 6000,
	ytick = {0,2000,4000,6000},
	ytick style = {font=\tiny},
    %legend image code/.code={\draw[#1] circle (0.1cm);}
]

\pgfplotstableread[col sep = comma]{data/stopdrinking-usagebymonth.csv}\mydata;
\addplot+[draw=ScienceDark,fill=none,thick] table[x=index,y=users] {\mydata};
\addplot+[draw=ArtsLight,fill=none,thick] table[x=index,y=threads] {\mydata};

\end{axis}
\end{tikzpicture}%
%
\begin{tikzpicture}[baseline]
\begin{axis}[IDEA line,
	% Graph Options
	no markers, % dense monthly data: the dash pattern tells the series apart
	width=.475*\columnwidth,	
	height=4cm,
	title=\texttt{r/Opiates\-Recovery},
	% X Axis Options
	xtick={1,13,25,37,49},                              %% Jim: apologies... this is a hack
    xticklabels={2014,2015,2016,2017,2018},             %%      will figure out what's going on here if we have time later
	%date coordinates in=x,
	%date ZERO=2014-01-01,
    %xticklabel=\year,
    %xmin=2014-01-01,
    %xmax=2018-01-01,
	xlabel = Year,
	%xtick ={1,2,3,4,5,6,7,8,9},
	% Y Axis Options
	ylabel = Distinct Users / Active Threads,
	ylabel style = {font=\tiny},
	ymin = 0, ymax = 800,
	ytick = {0,200,400,600,800},
	ytick style = {font=\tiny},
	legend style={at={(0.5,-1)},anchor=south}
    %legend image code/.code={\draw[#1] circle (0.1cm);}
]

\pgfplotstableread[col sep = comma]{data/opiaterecovery-usagebymonth.csv}\mydata;
\addplot+[draw=ScienceDark,fill=none,thick] table[x=index,y=users] {\mydata};
\addlegendentry{Distinct Users}
\addplot+[draw=ArtsLight,fill=none,thick] table[x=index,y=threads] {\mydata};
\addlegendentry{Active Threads}

\end{axis}
\end{tikzpicture}
%\setlength{\belowcaptionskip}{-20pt}
\caption{Distinct users and active threads for \texttt{r/stop\-drinking} (top) and \texttt{r/Opiates\-Recovery} (bottom) showing that both counts are experiencing an upwards trend over time for both subreddits.} 

\label{fig:activity}
\end{figure}



\newpage
%%%%%%%%%%%%%%%%%%%%%%%%%%%%%%%%%%%%%%%%%%%%%%%%%%%%%%%%%%
\subsection{Trends over Time: Stacked Bar Graph}
%%%%%%%%%%%%%%%%%%%%%%%%%%%%%%%%%%%%%%%%%%%%%%%%%%%%%%%%%%


\begin{figure*}[b!]
\Description{Five bar-chart panels stacked vertically and sharing an x axis of years from 1990 to 2015, one panel per type of paper, each bar a count of papers. Bibliographic: 20 in 1990, then between 0 and 4. Not Empirical: highest in 1992 (55) and 1998 (54), falling to single digits after 2010 except 31 in 2011. Explanatory: 3 in 1990, rising to a peak of 29 in 2006, then between 8 and 25. Design and Evaluation, the only panel in blue: between 10 and 39 with no clear trend. Descriptive: rising from 23 in 1990 to 58 in 2015.}
\begin{center}
\begin{tikzpicture}
\begin{groupplot}[
   group style={
       group size=1 by 5,
       x descriptions at=edge bottom,
       y descriptions at=edge left,
       vertical sep=1pt,
       horizontal sep=0pt},
  IDEA tufte panel, 
	% X Axis
	width=\columnwidth,
	xmin = 1990, xmax = 2015,
	xtick = {1990,1994,1998,2002,2006,2010,2015},
	x tick label style={font=\small,rotate=90, /pgf/number format/1000 sep=},
	axis line style={thick},
	% Y Axis
	ytick={10,20,30,40,50,60,70,80,90},
        ymin = 0,
        ymax = 50,
      ]


% Bibliographic
\nextgroupplot[bar width=7pt, height = 2.25cm, ymax = 25, title={Bibliographic}]
\addplot+[draw=none, fill=black!50] 
 coordinates {(1990,20) (1992,4) (1994,0) (1996,0) (1998,2) (2000,0) (2002,0) (2004,2) (2006,2) (2008,1) (2010,0) (2011,0) (2012,0) (2013,1) (2014,0) (2015,1) };       
  
% Not Empirical      
\nextgroupplot[bar width=7pt, height = 3.5cm, ymax=70, title={Not Empirical}]
\addplot+[draw=none, fill=black!50] 
coordinates {(1990,37) (1992,55) (1994,33) (1996,38) (1998,54) (2000,35) (2002,20) (2004,18) (2006,24) (2008,13) (2010,5) (2011,31) (2012,9) (2013,4) (2014,6) (2015,7) };

% Explanatory
\nextgroupplot[bar width=7pt, height = 3cm, title={Explanatory}]
\addplot+[draw=none, fill=black!50]  
coordinates {(1990,3) (1992,4) (1994,5) (1996,0) (1998,0) (2000,5) (2002,15) (2004,16) (2006,29) (2008,25) (2010,21) (2011,16) (2012,8) (2013,14) (2014,15) (2015,12) };

% D & E
\nextgroupplot[bar width=7pt, ymax=60, height = 3.5cm, title={Design and Evaluation}]
\addplot+[draw=none, fill=blue!50] 
coordinates {(1990,17) (1992,13) (1994,26) (1996,22) (1998,15) (2000,27) (2002,39) (2004,35) (2006,23) (2008,19) (2010,21) (2011,10) (2012,28) (2013,26) (2014,24) (2015,22) };

% Descriptive
\nextgroupplot[bar width=7pt, ymax = 70, height = 3.5cm, title={Descriptive}]
\addplot+[draw=none,fill=black!50]
coordinates {(1990,23) (1992,23) (1994,36) (1996, 40) (1998,29) (2000,32) (2002,27) (2004,29) (2006,23) (2008,42) (2010,53) (2011,42) (2012,55) (2013,55) (2014,55) (2015,58) };


\end{groupplot}
\end{tikzpicture}
\caption{An example of Tufte's stacked panel graph, which presents the same information as a stacked bar graph but better enables comparisons across each vertical dimension. Best practice is to use black and white when possible, with sparing colour for emphasis where appropriate. }
\label{stackedpanel}
\end{center}
\end{figure*}




% \begin{figure}[tb!]
% \begin{tikzpicture}
% \begin{groupplot}[
%   group style={
%       group size=1 by 5,
%       x descriptions at=edge bottom,
%       y descriptions at=edge left,
%       vertical sep=10pt,
%       horizontal sep=0pt},
%   IDEA line,
% 	% Graph Options
% 	height=3cm,
% 	 separate axis lines,
% 	% X Axis Options
% 	xmin=0, xmax= 35,
% 	xtick={1,5,10,15,20,25,30, 35},
% 	%date coordinates in=x,
% 	%date ZERO=2014-01-01,
%     %xticklabel=\year,
%     %xmin=2014-01-01,
%     %xmax=2018-01-01,
% 	xlabel = Number of NPC Requests,
% 	xlabel style = {font=\small\sffamily},
%     xticklabel style = {font=\small\sffamily},
% 	%xtick ={1,2,3,4,5,6,7,8,9},
% 	% Y Axis Options
% 	%ylabel = Proportion of Responses,
% 	ymajorgrids = false,
% 	ylabel shift = -.3cm,
% 	ymin = 0, ymax = 1,
% 	ytick = {0, 1.00},
% 	ylabel style = {font=\small\sffamily},
%     yticklabel style = {font=\small\sffamily},
% 	%ytick style = {font=\tiny},
% 	title style = {yshift=-.75cm,font=\sffamily\small}
% ]
% \pgfplotstableread[col sep = comma]{data/CBT_PropOverTime.csv}\cbtovertimedata;

% \nextgroupplot[xtick style = {white}, title style = {yshift=+.35cm,font=\sffamily\small}, title={Not Damaging}]
% \addplot+[mark=*, mark size=1.5pt, draw=gray, only marks,fill=gray,thin] table[x=Query,y=NotDamaging] {\cbtovertimedata};
% \addplot [solid, draw=black, fill=none,ultra thick] table[y={create col/linear regression={y=NotDamaging}}] {\cbtovertimedata};
% %\addlegendentry{Not Damaging}

% \nextgroupplot[xtick style = {white}, title={Empathize}]
% \addplot+[mark=*, mark size=1.5pt, draw=gray, only marks,fill=gray,thin] table[x=Query,y=Empathize] {\cbtovertimedata};
% \addplot [solid, draw=black, fill=none,ultra thick] table[y={create col/linear regression={y=Empathize}}] {\cbtovertimedata};
% %\addlegendentry{Empathize}

% \nextgroupplot[xtick style = {white}, title style = {yshift=+.25cm,font=\sffamily\small}, title={Reframe}]
% \addplot+[mark=*, mark size=1.5pt, draw=gray, only marks,fill=gray,thin] table[x=Query,y=Reframe] {\cbtovertimedata};
% \addplot [solid, draw=black, fill=none,ultra thick] table[y={create col/linear regression={y=Reframe}}] {\cbtovertimedata};
% %\addlegendentry{Reframe}

% \nextgroupplot[xtick style = {white}, title={Encourage}]
% \addplot+[mark=*, mark size=1.5pt, draw=gray, only marks,fill=gray,thin] table[x=Query,y=Encourage] {\cbtovertimedata};
% \addplot [solid, draw=black, fill=none,ultra thick] table[y={create col/linear regression={y=Encourage}}] {\cbtovertimedata};
% %\addlegendentry{Encourage}

% \nextgroupplot[title={Solution}]
% \addplot+[mark=*, mark size=1.5pt, draw=gray, only marks,fill=gray,thin] table[x=Query,y=Solution] {\cbtovertimedata};
% \addplot [solid, draw=black, fill=none,ultra thick] table[y={create col/linear regression={y=Solution}}] {\cbtovertimedata};
% %\addlegendentry{Solution}


% %\addplot+[mark=x, only marks,draw=Monster_0,fill=none,semithick] table[x=Query,y=Proportion, y error plus=CI-upper, y error minus=CI-lower] {\barbariansolutiondata};
% %\addlegendentry{Monster}
% %\addplot [draw=Monster_0, fill=none,ultra thick] table[y={create col/linear regression={y=Proportion}}] {\barbariansolutiondata};

% \end{groupplot}
% \end{tikzpicture}

% %\setlength{\belowcaptionskip}{-20pt}
% \caption{The proportion of responses that were classified as meeting five criteria: Empathize, Reframe, Encourage, Solution, and Not Damaging. The three categories related to CBT, Empathize, Reframe, and Encourage, were found to decrease over time ($p=0.05$). However, we did not find these changes for Solution and Not Damaging. 
% %decreased over time for all participants, however responses by participants in the \textsc{Cleric} group decreased more quickly than those in the \textsc{Monster} group.
% } 

% \label{fig:SolutionsOverQueryNumber}
% \end{figure}



% \newpage
% %%%%%%%%%%%%%%%%%%%%%%%%%%%%%%%%%%%%%%%%%%%%%%%%%%%%%%%%%%
% \subsection{Illustrating Scenarios: Rotograph}
% %%%%%%%%%%%%%%%%%%%%%%%%%%%%%%%%%%%%%%%%%%%%%%%%%%%%%%%%%%

% - You don't need to be an artist
% - take a picture, trace over it
% - often by far the most effective technique for simply demonstrating key concepts in your paper 




\newpage
%%%%%%%%%%%%%%%%%%%%%%%%%%%%%%%%%%%%%%%%%%%%%%%%%%%%%%%%%%
\subsection{PRISMA Flow Diagram}
%%%%%%%%%%%%%%%%%%%%%%%%%%%%%%%%%%%%%%%%%%%%%%%%%%%%%%%%%%

- Example provided by Katja Rogers
- http://www.prisma-statement.org
- ``PRISMA is an evidence-based minimum set of items for reporting in systematic reviews and meta-analyses. PRISMA focuses on the reporting of reviews evaluating randomized trials, but can also be used as a basis for reporting systematic reviews of other types of research, particularly evaluations of interventions.''


\pgfdeclarelayer{background}
\pgfsetlayers{background,main}
\begin{figure}[b!]
\Description{PRISMA flow diagram with four stages labelled down the left side: Identification, Screening, Eligibility and Included. Records identified through database searching (N=1035) lead to records after duplicates removed (N=1030), then records screened (N=1030), of which 239 records were excluded. 791 full-text articles were assessed for eligibility, of which 731 were excluded with reasons. 60 studies were included in qualitative and quantitative synthesis.}
{
\footnotesize
\begin{tikzpicture}
    \tikzstyle{block}=[
        draw,
        thick,
        text width=5cm,
        minimum height=1cm,
        align=center
    ]
    \tikzstyle{header}=[rotate=90, align=center]
    \tikzstyle{line}=[-latex, draw, very thick]
    \node[block] (included-0) {Studies included in qualitative and quantitative synthesis (N=60)};
    \node[block, above=2em of included-0] (eligibility-1) {Full-text articles assessed for eligibility \\ (N=791)};
    \node[block, right=2em of eligibility-1] (eligibility-2) {Full-text articles excluded, with reasons \\ (N=731)};
    \node[block, above=2em of eligibility-1] (screening-1) {Records screened \\ (N=1030)};
    \node[block, right=2em of screening-1] (screening-2) {Records excluded \\ (N=239)};
    \node[block, above=2em of screening-1] (screening-0) {Records after duplicates removed\\ (N=1030)};
    \node[block, above=2em of screening-0, minimum height=1.5cm] (identification-1) {Records identified through database searching \\ (N=1035)}; 
    \coordinate (screening-mid) at ($ (screening-0.west)!0.5!(screening-1.south west) $);
    \coordinate (identification-mid) at ($ (screening-0.west)!0.5!(identification-1.north west-|screening-0.west) $);
    \node[header, left=1cm of included-0.west, anchor=north] (included) {\textbf{Included}}; %left=\textwidth/6 of included-0.west, anchor=north
    \node[header, left=1cm of eligibility-1.west, anchor=north] (eligibility) {\textbf{Eligibility}};
    \node[header, left=1cm of screening-1.west, anchor=north] (screening) {\textbf{Screening}};
    \node[header, left=1cm of identification-mid, anchor=north] (identification) {\textbf{Identification}};
    % Draw edges
    \path [line] (identification-1) -- ([yshift=.5cm]screening-0.north) -- (screening-0);
    \path [line] (screening-0) -- (screening-1);
    \path [line] (screening-1) -- (screening-2);
    \path [line] (screening-1) -- (eligibility-1);
    \path [line] (eligibility-1) -- (eligibility-2);
    \path [line] (eligibility-1) -- (included-0);
\end{tikzpicture}
}
\caption{Number of results at each stage of the review, as represented in a PRISMA flow diagram.}
\end{figure}




\newpage
%%%%%%%%%%%%%%%%%%%%%%%%%%%%%%%%%%%%%%%%%%%%%%%%%%%%%%%%%%
\subsection{Scheduling: Gantt Charts}
%%%%%%%%%%%%%%%%%%%%%%%%%%%%%%%%%%%%%%%%%%%%%%%%%%%%%%%%%%

- supported by the \texttt{pgfgantt} package

\begin{figure}[b!]
\Description{Gantt chart of a research plan from May 2018 to April 2023 under five yearly headings, Year 1: 2019 -- 2020 to Year 5: 2023 -- 2024. Objective 1, Design Framework (May 2018 to April 2021): literature review, artifact review, framework development, then mentoring and publication. Objective 2, Interaction Mechanics (May 2020 to April 2022): design, development, lab study, then analysis and publication. Objective 3, In-Situ Evaluation (May 2021 to April 2023): requirements gathering, development, ecological validation, then analysis and publication. Within each objective the activities follow one another with small overlaps, and the objectives overlap by about a year.}
\begin{ganttchart}[
    x unit=.208cm,
    y unit title=0.4cm,
    y unit chart=0.3cm,
    vgrid,
    time slot format=isodate-yearmonth,
    time slot unit=month,
    vgrid={*{3}{draw=none},dotted,*{3}{draw=none},dotted,*{3}{draw=none},*{1}{draw=gray}},
    %
    % TITLE OPTIONS
    %
    title/.style={draw=none, fill=none},
    title label node/.append style={below=-1.6ex},
    title top shift=-.5,
    title label font=\tiny,
    %
    % BAR OPTIONS
    %
    bar height=.5,
    bar/.append style={draw = none, fill=blue!50},
    bar label font=\tiny,
    %   
    % GROUP OPTIONS
    %
    group label font={\tiny},
    group right shift=0,
    group top shift=.6,
    group height=.3,
    group peaks height=.2,
   ]{2018-05}{2023-04}

    
    \gantttitle{\textbf{Year 1: 2019 -- 2020}}{12}
    \gantttitle{\textbf{Year 2: 2020 -- 2021}}{12}
    \gantttitle{\textbf{Year 3: 2021 -- 2022}}{12}
    \gantttitle{\textbf{Year 4: 2022 -- 2023}}{12}
    \gantttitle{\textbf{Year 5: 2023 -- 2024}}{12}
                
\ganttset{progress label text={}, link/.style={black, -to}}

\ganttgroup{\textbf{Objective 1: Design Framework}}{2018-05}{2021-04} \\
\ganttbar{Literature Review (1.1)}{2018-05}{2019-03} \\
\ganttbar{Artifact Review (1.2)}{2018-09}{2019-07} \\
\ganttbar{Framework Development (1.3)}{2019-05}{2020-04} \\
\ganttbar{Mentoring \& Publication (1.4)}{2020-05}{2021-04} \\

\ganttgroup{\textbf{Objective 2: Interaction Mechanics}}{2020-05}{2022-04} \\
\ganttbar{Design (2.1)}{2020-05}{2020-07} \\
\ganttbar{Development (2.2)}{2020-08}{2020-12} \\
\ganttbar{Lab Study (2.3)}{2021-01}{2021-07} \\
\ganttbar{Analysis \& Publication (2.4)}{2021-07}{2022-04} \\

\ganttgroup{\textbf{Objective 3: In-Situ Evaluation}}{2021-05}{2023-04} \\
\ganttbar{Requirements Gathering (3.1)}{2021-05}{2021-10} \\
\ganttbar{Development (3.2)}{2021-11}{2022-06} \\
\ganttbar{Ecological Validation (3.3)}{2022-06}{2022-12} \\
\ganttbar{Analysis \& Publication (3.4)}{2023-01}{2023-04} 


\end{ganttchart}
\caption{Research activities are designed to scaffold development of a design framework and evaluation methodology, culminating in an in-situ evaluation with Homewood Research Institute.}
\label{gantt}
\end{figure}


\section{Referencing}
\label{sec:referencing}

\begin{markdown}


# Reporting Measurements and Statistics
`\label{sec:statistics}`{=tex}

As part of loading this template, the `siunit` package is loaded, which supports some common requests, like metric units and angles. This template also includes some macros that help with reporting of common elements of HCI papers like display measurements, study designs, and statistical tests. See below for some common examples: 

## Sizes

The `siunitx` package already has a feature for angles, `\ang{45}`{=tex}.

The display resolution was `\size[pixels]{1920}{1080}`{=tex}.
The volume of the cube was `\sizen[\cm]{13;12;5}`{=tex}.

We performed a 5 (width)`\by`{=tex} 2 (distance)`\by`{=tex} 2 (angle) factorial ANOVA.


## Summary Statistics

We recruited 20 participants (`\ages{23}{51}`{=tex}; `\gender{12}{women}`{=tex}, `\gender{7}{men}`{=tex}, and `\gender{1}{transgender}`{=tex}).
Alternatively: `\genderlist{12;7;1}{women, men, and transgender}`{=tex}.

A mean and standard deviation (`\mean[\cm]{32.1}`{=tex}, `\sd{4.1}`{=tex}) with same units (`\msd[pixels]{3.3333333}{0.32}`{=tex}).



## Reporting Statistical Tests

### $p$-values

`\p{.03}`{=tex}

### ANOVA & $F$-statistics

Some made up stats (`\F{2}{3}{4.33}`{=tex}, `\p{.03}`{=tex}, `\etasqp{.67}`{=tex})
Some made up stats (`\F{2}{3}{15.33}`{=tex}, `\p{.003}`{=tex}, `\etasqp{.67}`{=tex}).
Some made up stats (`\anovae{2}{3}{4.33}{.03}{.6666667}`{=tex})
Some made up stats (`\anovae{2}{3}{15.333333333}{.003}{.6666667}`{=tex})
Some made up stats (`\anovae{2}{3}{30.55555}{.0000000055}{0.9666667}`{=tex})
Effect sizes always print as a number, even when small (`\anovae{1}{38}{0.47}{.497}{.012}`{=tex})

### CHI-Square Test

Here is a chi-square test (`\chisquare{33}{3.323}{.000005}`{=tex}).


\end{markdown}

%%%%%%%%%%%%%%%%%%%%%%%%%%%%%%%%%%%%%%%%%%%%%%%%%%%%%%%%%%%%%%%%%%%%%%%%%%%%%
%\section{Limitations}
%%%%%%%%%%%%%%%%%%%%%%%%%%%%%%%%%%%%%%%%%%%%%%%%%%%%%%%%%%%%%%%%%%%%%%%%%%%%%



% https://transparenthci.medium.com/a-reviewer-asked-you-to-make-your-materials-open-now-what-f4af38e44b86
\newpage
\section{Transparency} 
\label{sec:transparency}

- OSF

%%%%%%%%%%%%%%%%%%%%%%%%%%%%%%%%%%%%%%%%%%%%%%%%%%%%%%%%%%%%%%%%%%%%%%%%%%%%%
%\section{Conclusions}
%%%%%%%%%%%%%%%%%%%%%%%%%%%%%%%%%%%%%%%%%%%%%%%%%%%%%%%%%%%%%%%%%%%%%%%%%%%%%




%%%%%%%%%%%%%%%%%%%%%%%%%%%%%%%%%%%%%%%%%%%%%%%%%%%%%%%%%%%%%%%%%%%%%%%%%%%%%
%\section{Acknowledgements}
%%%%%%%%%%%%%%%%%%%%%%%%%%%%%%%%%%%%%%%%%%%%%%%%%%%%%%%%%%%%%%%%%%%%%%%%%%%%%

\begin{markdown}

# Acknowledgements

This template is a compilation of common mistakes, reviewer feedback, and paper writing experience. It is intended to help identify the ‘little things’ that often get in the way of a successful CHI submission, and may be helpful in reviewing your paper `\textit{well before}`{=tex} the deadline. For example, to guide peer reviews, such as reviewing circles, or to help communicate necessary revisions. 

While I compiled the different examples, tips, and hacks in one place, they incorporate the experience, advice, and feedback from many others I've worked with, including (in alphabetical order): Robert Gauthier, Mark Hancock, Lennart Nacke, Adrian Reetz, Katja Rogers, Oliver Schneider, Stacey Scott, and Dan Vogel. Suggestions are always welcome, please email to `\href{mailto:james.wallace@uwaterloo.ca}{james.wallace@uwaterloo.ca}`{=tex}. 

\end{markdown}


\printbibliography


%\input{submission/11-appendix.tex}

\newpage

%\pagestyle{fancy}
\rhead{\href{http://hci.uwaterloo.ca}{\includegraphics[width=3cm]{figures/uwhci-full-cs.png}}}
\lhead{\textcolor{gray}{\tiny{v3.0 -- \today}}}
%\renewcommand{\headrulewidth}{0pt}

\newenvironment{checklist}{%
  \begin{list}{}{}% whatever you want the list to be
  \let\olditem\item
  \renewcommand\item{\olditem[$\Box$] }
}{%
  \end{list}
}

\appendix
%%%%%%%%%%%%%%%%%%%%%%%%%%%%%%%%%%%%%%%%%%%%%%%%%%%%%%
\section{\textsc{CHI Paper Checklist}}
%%%%%%%%%%%%%%%%%%%%%%%%%%%%%%%%%%%%%%%%%%%%%%%%%%%%%%
\label{sec:checklist}

\noindent\makebox[\linewidth]{\rule{\linewidth}{0.4pt}}
% %%%%%%%%%%%%%%%%%%%%%%%%%%%%%%%%%%%%%%%%%%%%%%%%%%%%%%
% \subsection*{\textsc{The Big Picture}}
%%%%%%%%%%%%%%%%%%%%%%%%%%%%%%%%%%%%%%%%%%%%%%%%%%%%%%
%\vspace{0.5cm}

You should answer the following in \emph{each of} the abstract, introduction, and related work sections to varying levels of detail. For instance, you might answer each question using about 1 sentence in the abstract, 1 paragaph in the introduction, and 1--3 paragraphs in the related work. If your reviewers can't answer these questions after reading the introduction of your paper, the odds that they can argue for accepting it are reduced drastically. 

\begin{checklist}
    \item What problem are you trying to solve?
    \item Who cares if you do? 
    \item What's the solution to the problem? 
	\item What is the specific contribution of your paper? (e.g., Precisely how does your design build on existing work? or What new findings does this study provide over similar existing studies?)
	\item Have you summarized your findings? Are take-home messages clear and meaningful? 
\end{checklist}

When thinking about how to articulate your contribution, consider framing it in terms of \citet{wobbrock2016research}, who suggest seven HCI contributions: Empirical, Artifact, Methodological, Theoretical, Dataset, Survey, and Opinion.

%\item Is the research described clearly? Could a fellow student replicate the work if all they had was this paper?  
	

\noindent\makebox[\linewidth]{\rule{\linewidth}{0.4pt}}
%%%%%%%%%%%%%%%%%%%%%%%%%%%%%%%%%%%%%%%%%%%%%%%%%%%%%%
\section*{\textsc{Data Collection, Analysis, and Results}}
%%%%%%%%%%%%%%%%%%%%%%%%%%%%%%%%%%%%%%%%%%%%%%%%%%%%%%

Be sure to report and describe: 
\begin{checklist}
	\item What data was collected? (e.g., independent variables) And how? (e.g., computer logs, video recordings, field notes, questionnaires, etc.) 

	\item What qualitative data analysis method(s) were conducted (e.g., high-level video review, affinity diagramming on your interview data, in-depth video coding, etc.) and by who(m)?
	
	\item What quantitative data analysis method(s) were used? (e.g., One-Way Analysis of Variance (ANOVA), Repeated-measures Analysis of Variance (RM-ANOVA), t-test, regression, etc.) 	
	
	\begin{enumerate}
 		\item State your alpha value, and any transformations or adjustments that were applied (e.g., for family-wise errors)
		\item Report summary statistics of your actual data (e.g. mean, standard deviations, or standard error) as well as the statistical test results ($t$-values, $F$-values, $p$-values, effect size, etc.)
	\end{enumerate}
	\item When including any figures or tables of your data, always indicate directly in the figure or table any statistically significant differences between data values
\end{checklist}


\noindent\makebox[\linewidth]{\rule{\linewidth}{0.4pt}}
%%%%%%%%%%%%%%%%%%%%%%%%%%%%%%%%%%%%%%%%%%%%%%%%%%%%%%
\section*{\textsc{Figures and Tables}}
%%%%%%%%%%%%%%%%%%%%%%%%%%%%%%%%%%%%%%%%%%%%%%%%%%%%%%
Do all of the Figures and Tables:
\begin{checklist}
	\item Clearly highlight the relevant information? Can they be (mostly) understood without reading the paper? 
	\item Have a caption and title? Are they annotated? 
	\item Have labelled axes, and appropriate units / scale?
	\item Are they legible when printed? In black and white?
	\item Are they referenced in the body text? Are they numbered sequentially?
\end{checklist}











%\noindent\makebox[\linewidth]{\rule{\linewidth}{0.4pt}}
%%%%%%%%%%%%%%%%%%%%%%%%%%%%%%%%%%%%%%%%%%%%%%%%%%%%%%
\section*{\textsc{References}}
%%%%%%%%%%%%%%%%%%%%%%%%%%%%%%%%%%%%%%%%%%%%%%%%%%%%%%

\begin{checklist}
	\item Have you \begin{em}\textbf{thoroughly}\end{em} checked your reference list at the end of your paper?
	\item Have you checked references to author's names in the main text for spelling?
\end{checklist}

\begin{center}

\small % The size of the table text can be changed depending on content. Remove if desired.
\def\arraystretch{2}
\begin{tabular}{p{6cm} p{4cm} p{4cm}}
& \textbf{YES} & \textbf{NO} \\
 \cmidrule{2-2}  \cmidrule{3-3} 
Include in-line citations in sentence & your claim [1]. & your claim. [1]. \\
Citations follow author's name & Bauer et al. [1] claimed X, Y, and Z. & Bauer et al. claimed X, Y, and Z [1]. \\
Citations are not nouns & ... as Bauer et al. [1] found ... & ... as [1] found ... \\
Citations immediately follow claims & Previous research has found X [1], Y [2,3], and Z [4, 5]. &  Previous research has found X, Y, and Z [1, 2, 3, 4, 5]. \\
All claims backed by a citation & The long-term use of smartphones causes increased anxiety [1, 2, 4]. & The long-term use of smartphones causes anxiety.  \\
Summarize relevant citations, including findings from previous studies & Bauer et al. [1] previously showed that FitBit's have long-term benefits for mental health. They found initial positive health benefits ... & Bauer et al. [1] previously studied the impact of FitBit on long-term mental health. \\
\end{tabular}
\end{center}


% \noindent\makebox[\linewidth]{\rule{\linewidth}{0.4pt}}
% %%%%%%%%%%%%%%%%%%%%%%%%%%%%%%%%%%%%%%%%%%%%%%%%%%%%%%
% \section*{\textsc{Formatting \& Misc.}}
% %%%%%%%%%%%%%%%%%%%%%%%%%%%%%%%%%%%%%%%%%%%%%%%%%%%%%%
% \vspace{0.5cm}

%\begin{checklist}
% 	\item Did you use the most recent \href{https://chi2027.acm.org/chi-publication-formats/}{ACM SIGCHI Paper Format}?
% 	\item Are all acronyms defined on first use? 
% \end{checklist}

% \begin{checklist}
% 	\item When writing about numbers less than or equal to ten, use full words: ``one'', ``two'' ...etc, 
% 	\item When a sentence starts with a number, it should also use words (e.g., ``Twenty-four participants were recruited...'') 
% 	\item Be consistent (e.g., ``There were 23 participants who wrote CHI papers carefully, and only 2 that wrote them at the last minute'')
% \end{checklist}



\noindent\makebox[\linewidth]{\rule{\linewidth}{0.4pt}}
%%%%%%%%%%%%%%%%%%%%%%%%%%%%%%%%%%%%%%%%%%%%%%%%%%%%%%
\section*{\textsc{Accessibility}}
%%%%%%%%%%%%%%%%%%%%%%%%%%%%%%%%%%%%%%%%%%%%%%%%%%%%%%
From the SIGCHI \href{https://sigchi.org/resources/guides-for-authors/accessibility/}{Accessibility Guide for Authors}, which CHI asks authors to follow (see \autoref{sec:accessibility}):
\begin{checklist}
    \item Does every figure have a \texttt{\textbackslash Description\{...\}}? Does it describe what the figure shows (chart type, axes, key values or trends) rather than repeat the caption? Is there no \texttt{CHIZen Warning} about a missing description in the log?
    \item Can every figure be read without colour? (``Don't rely only on colour'': use shapes, patterns or direct labels as well.)
    \item Is every table a real table, not an image? Are its header cells marked (the first row by default; otherwise set \texttt{table/header-rows} with \texttt{\textbackslash tagpdfsetup}), and does every row have the same number of cells?
    \item Is every equation marked up as maths, not an image?
    \item Are headings and lists made with \texttt{\textbackslash section} and \texttt{itemize}/\texttt{enumerate}, not bold text and manual numbers?
    \item Is the document metadata right: the real title in \texttt{\textbackslash title}, and the language in \texttt{\textbackslash DocumentMetadata} (\texttt{lang = en})?
    \item Does link text say where the link goes (not ``here'')?
    \item Did you compile with LuaLaTeX and check the PDF with a validator (veraPDF, or Acrobat's accessibility checker)?
    \item For the camera-ready version: did you also enter the figure descriptions in PCS?
\end{checklist}


\noindent\makebox[\linewidth]{\rule{\linewidth}{0.4pt}}
%%%%%%%%%%%%%%%%%%%%%%%%%%%%%%%%%%%%%%%%%%%%%%%%%%%%%%
\section*{\textsc{Before You Submit }}
%%%%%%%%%%%%%%%%%%%%%%%%%%%%%%%%%%%%%%%%%%%%%%%%%%%%%%
\begin{checklist}
    \item \textbf{1 week before deadline} Have you checked your paper for everything above? 
    \item \textbf{1 week before} Have you asked a friend to check your paper for everything above? 
    \item \textbf{1 day before} Check with your co-authors that everything is ready
    \item \textbf{After the deadline} Email a copy of the submitted PDF to your co-authors
    \item \textbf{Camera-ready} Have you switched the paper to final mode (\texttt{\textbackslash def\textbackslash chizenmode\{final\}}), with up-to-date figure PDFs in \texttt{figures/tikz/}, so that it loads only packages ACM's TAPS accepts? (See the manual, Section ``Camera-Ready for ACM TAPS''.)
\end{checklist}





%\noindent\makebox[\linewidth]{\rule{\linewidth}{0.4pt}}
%%%%%%%%%%%%%%%%%%%%%%%%%%%%%%%%%%%%%%%%%%%%%%%%%%%%%%
\section*{\textsc{BONUS: Words and Phrases to Avoid}}
%%%%%%%%%%%%%%%%%%%%%%%%%%%%%%%%%%%%%%%%%%%%%%%%%%%%%%
\begin{multicols}{5}
\small
\begin{itemize}
\setlength{\itemsep}{0.5pt}
\item[] actually
\item[] almost
\item[] always
\item[] anyway
\item[] basically
\item[] best
\item[] believe
\item[] can't
\item[] clearly
\item[] don't
\item[] essentially
\item[] feel
\item[] folks
\item[] good
\item[] guy
\item[] hassle
\item[] I
\item[] impossible
\item[] kids
\item[] kind of
\item[] like
\item[] lots
\item[] many
\item[] never
\item[] nice
\item[] obviously
\item[] OK
\item[] of course
\item[] pretty good
\item[] prove(n)
\item[] quite
\item[] really
\item[] so
\item[] stuff
\item[] thing
\item[] totally
\item[] user
\item[] utilize
\item[] various
\item[] very
\item[] well
\item[] won't
\item[] wouldn't
\item[] you
\end{itemize}

\end{multicols}


%\item Has a friend proof-read the full draft of your paper? 
	






\end{document}
\endinput
%%
%% End of file.
